% !TEX root = ../main.tex
\chapter{Research Methodology}
\label{ch:methodology}

Chapter~\ref{ch:literature} ended with five gaps. This chapter gives each of them a measurable procedure and sets out how EndorseRank, Adaptive Weighted PageRank (AWP) and the coupled walks built from their two layers are evaluated as reputation scores for smart contracts on Arbitrum One. EndorseRank runs on the ERC-20 allowance graph, and AWP \parencite{do2023awp} on the transfer graph. Both weight their edges as AWP's authors specify and are solved with weighted PageRank \parencite{page1999pagerank}; AWP restarts the walk in proportion to sending activity, and EndorseRank restarts it uniformly. Every scored and evaluated address is a contract. Externally owned accounts (EOAs) enter the graphs as owners, senders and recipients and are never scored. Both graphs keep the approvals and transfers of five widely traded tokens and value every amount in US dollars (Section~\ref{sec:amount-revision}).

The design follows construct alignment. Reputation scores are continuous rankings, and the checks rest on observable contract-level statistics instead of binary default labels. Each ranking method is treated as one indicator of a reputation construct that cannot be observed directly, in the sense of construct validity \parencite{cronbach1955construct}. The claim is comparative, every score set against the checks of each construct and against the raw degrees it smooths; it is not a prediction of trading profit, of credit default or of the safety of a contract's code.

Every procedure runs on public blockchain logs and public daily token prices inside fixed windows. Seeds are deterministic and the evaluation code is public, so anyone who reruns the extraction, or holds the downloaded pair tables, can regenerate the reported tables.

Four principles guide each operational choice.
\begin{itemize}
\item Same-seed comparability: every score is computed on the graphs of one seed set, the matched contract cohort, and evaluated on identical sets of contracts, so no difference between scores can come from sampling.
\item Fixing choices in advance: the plan wrote down every choice that decides a result before the data it is judged on, and the revision of the amounts was written down before any revised quantity was computed, though after the registered results were known. The events took place in this order.\footnote{All dates and times are given in UTC. The commits record the local time of the machine on which they were made (UTC+9), in which the registration commit carries the date 10~October~2026.}
\begin{enumerate}
\item At 13:41~UTC on 9~October~2026 the analysis plan and the versioned configuration were committed to the public repository (commit \texttt{43d8045}) before any log of the study was extracted. They fix the windows, the cohort rule, the account-type rule, every score with its damping, decay, value transform and restart rule, the values of $\lambda$, the proxies, the labels of the holdout W0, the bootstrap resamples and seed, the benchmark protocol, the robustness checks and rule~A for the coupled operator (Section~\ref{sec:holdout-protocol}).
\item The same day, scan~A extracted the observation window. The cohort was fixed, and the ego events of the cohort, the pair tables, the proxies and the W0 labels were built.
\item The same-window and W0 results were computed, and rule~A was judged.
\item At 16:49~UTC on the same day the registration of the window W1 was committed (commit \texttt{0d3a057}): rule~B with its margin, three further results reported in a fixed order, two sensitivity analyses, and the comparison on outcomes built from neither edge type with the wording of its conclusions (Sections~\ref{sub:fresh-holdout} and~\ref{sub:neutral-label-protocol}). The script that runs scan~B refuses to run without this commit and records the hashes of the registered files.
\item Scan~B extracted July to September 2026, the W1 labels were counted, and the registered evaluation was run.
\item After the W1 labels were known came the supplementary contrasts (EndorseRank with activity restarts against the in-approve degree, and EndorseRank against AWP in W1), the robustness checks, the Sybil model and the benchmark.
\item On 10~October~2026, after these results were known, the amounts were revised (Section~\ref{sec:amount-revision}). Every amount is valued in US dollars, only five tokens with reliable prices are kept, the value transform takes its slope from the median amount, and the cohort, the graphs, the proxies and the spender labels are rebuilt on those tokens. The revision was committed between 07:54 and 11:24~UTC, before any revised quantity was computed, and the revised analysis was then run with the protocols, rules and thresholds of the registered one. Chapter~\ref{ch:results} reports the revised analysis.
\end{enumerate}
In the registered analysis, only rule~A (its two contrasts in W0 and its same-window part) and the results registered for W1 (rule~B, the three further results, the two sensitivity analyses and the comparison on outcomes built from neither edge type) were fixed before the data they are judged on. In the revised analysis none was: the same rules and thresholds are applied, but to scores defined after the registered results were known. Every other contrast, the contrast of each PageRank with its own raw degree among them, was computed after the data it is judged on were known; it is reported as such, with its interval.
\item Construction and timing: the pair tables are aggregated once for each anchor date and loaded before any timer starts, and the benchmark times the solver call.
\item Scope: H1 and Claim~C3 compare the orderings of EndorseRank and AWP; H2, H3 and Claims C1, C2 and C4 concern every score, C5 the coupled operator and its ablation, and C6 the approval and transfer counts and their PageRanks on trader outcomes. RQ4 and RQ5 draw on the spender cohorts of W0 and W1; rule~B and the comparison on outcomes built from neither edge type add the trader cohort.
\end{itemize}

\dissertationSubheading[sec:notation]{Notation}

The notation used in this chapter and in Chapters~\ref{ch:results} and~\ref{ch:discussion} is collected in Table~\ref{tab:notation-methods} below, and a condensed list appears in Table~\ref{tab:notation}. Vectors are set in bold and sets in calligraphic type. Addresses are treated as opaque identifiers; we do not map them to any external identity.

\begingroup
\setstretch{1}\small\setlength{\tabcolsep}{3.5pt}\captionsetup{font=normalsize}
\begin{longtable}{lp{0.70\linewidth}}
\caption{Notation for graphs, scores, proxies, and parameters.}
\label{tab:notation-methods}\\
\toprule
Symbol & Meaning \\
\midrule
\endfirsthead
\caption[]{(continued)}\\
\toprule
Symbol & Meaning \\
\midrule
\endhead
\bottomrule
\endlastfoot
$\mathcal{W}$ & Matched contract cohort ($|\mathcal{W}|=n=15{,}107$), the seed set of both graphs \\
$\mathcal{C}$ & Contracts evaluated in one analysis: $\mathcal{W}$, or a spender or trader cohort of a holdout \\
$G=(V,E)$ & Graph passed to the PageRank solver: $V$ holds the end points of the retained edges $E$ \\
$G_{\mathrm{approve}}=(V_A,E_A,w_A)$ & EndorseRank allowance graph \\
$G_{\mathrm{transfer}}=(V_T,E_T,w_T)$ & AWP transfer graph \\
$a^{\star}(u,v,\mathrm{token})$ & Latest allowance of owner $u$ to spender $v$ on one token up to the anchor (US dollars at the anchor-day price) \\
$\Delta t^{\star}(u,v,\mathrm{token})$ & Age of the approval that set $a^{\star}$ \\
$w_A(u,v)$ & EndorseRank edge weight, owner $u\to$ spender $v$: $\sigma(\Delta t^{\star})V(a^{\star})$ summed over tokens \\
$w_T(u,v)$ & AWP edge weight, sender $u\to$ receiver $v$: $\sigma(\Delta t_e)V(x_e)$ summed over transfers \\
$x_e$ & Amount of transfer $e$ in US dollars at the price of its day \\
$V(z)$, $b$ & Bounded value transform of AWP, $V(z)=2/(1+e^{-bz})-1$, with $b=\ln 3/m$ on US-dollar amounts ($b_T$ for transfers, $b_A$ for allowances) \\
$m_T$, $m_A$ & Median US-dollar amounts of the transfers and of the finite latest allowances up to $t_1$ \\
$U$ & Largest amount at which an allowance counts in the allowance layer of C-PR \\
$X_u$ & Activeness of address $u$, the restart weight of AWP \\
$\tilde w_A(u,v)$, $\tilde w_T(u,v)$ & Weights of the two layers of C-PR, in US dollars \\
$o(u)$; $o_A(u)$, $o_T(u)$ & Out-strength $\sum_v w(u,v)$ of node $u$; the same in the allowance and transfer layers \\
$P$ & Transition matrix, $P_{u,v}=w(u,v)/o(u)$; rows of dangling nodes are zero \\
$\mathbf{1}[\cdot]$ & Indicator: $1$ if the condition holds, $0$ otherwise \\
$\mathbf{1}_{\mathrm{dang}}$ & Indicator vector of the dangling nodes, entries $\mathbf{1}[o(u)=0]$ \\
$\mathbf{r}$, $\mathbf{r}^{(t)}$ & PageRank vector and its iterate after $t$ steps \\
$\boldsymbol{\pi}$ & Restart (teleportation) distribution, also used for dangling mass: uniform for EndorseRank and C-PR, proportional to $X_u$ for AWP, and for S-PR the allowance layer's score restricted to $V_T$ and renormalized \\
$d$ & PageRank damping factor (default $0.85$) \\
$\varepsilon$ & PageRank convergence tolerance ($10^{-8}$) \\
$I_{\max}$ & Maximum PageRank iterations ($300$) \\
$\lambda$ & Allowance-layer weight of the coupled operator C-PR ($\lambda=0.5$ primary; $0.25$, $0.75$ sensitivity; at $\lambda=1$ and $\lambda=0$ one layer is walked alone) \\
$\alpha_A(u)$, $\alpha_T(u)$ & Layer weights in node $u$'s C-PR row; $D_{\lambda}(u)$ is their common denominator \\
$P^{(\lambda)}$ & Coupled transition matrix of C-PR \\
$\mathbf{s}_M$ & Reputation score vector for method $M$ \\
$\mathbf{p}_q$ & Proxy value vector for proxy $q$ \\
$\sigma(\Delta t)$ & Logistic time-decay on the age of an approval or transfer \\
$k$, $t_0$ & Decay steepness ($0.01$) and midpoint ($180$ days) \\
$\Delta t$ & Age of an event in days before the decay anchor ($T_{\mathrm{obs}}$; $t_1$ in W0) \\
$T_{\mathrm{obs}}$ & End of the observation window and freeze of W1 (2026-06-30 23:59:59 UTC) \\
$t_1$, $t_2$ & Score freeze of W0 (2026-03-31 23:59:59 UTC) and end of its label window (2026-06-30 23:59:59 UTC) \\
W0, W1 & Holdout windows: labels from April to June 2026 with scores at $t_1$, and from July to September 2026 with scores at $T_{\mathrm{obs}}$ \\
$\delta$ & Non-inferiority margin of rule~B ($0.01$) \\
$B$ & Paired bootstrap resamples ($400$; $2{,}000$ in Section~\ref{sub:neutral-label-protocol}) \\
$\Delta\tau$ & Paired difference of two Kendall coefficients on the same resamples \\
$\rho$ & Spearman's rank correlation \\
$R_i$, $S_i$ & Mid-ranks of contract $i$'s score and proxy value (tied values share their average position) \\
$\tau$, $\tau_b$ & Kendall's rank correlation; every reported value is the tie-corrected $\tau_b$ \\
$n_c$, $n_d$ & Numbers of concordant and discordant contract pairs \\
$n_0$; $n_1$, $n_2$ & All pairs, $n(n-1)/2$; pairs tied in the first and in the second ranking \\
\end{longtable}
\endgroup

Except where stated otherwise, PageRank means the weighted variant of the algorithm of \textcite{page1999pagerank}, whose update at iteration $t$ is
\begin{equation}
\label{eq:pagerank}
\mathbf{r}^{(t+1)}
=
(1-d)\,\boldsymbol{\pi}
+
d\,P^{\top}\mathbf{r}^{(t)}
+
d\,\bigl(\mathbf{r}^{(t)\top}\mathbf{1}_{\mathrm{dang}}\bigr)\,\boldsymbol{\pi},
\end{equation}
where $P$ is the transition matrix of Equation~\eqref{eq:transition-impl} below, $\boldsymbol{\pi}$ is the restart distribution over the nodes, uniform ($\pi_u=1/|V|$) unless a method specifies another (AWP in Section~\ref{sub:awp-graph}, S-PR in Section~\ref{sub:coupled-graph}), and $\mathbf{1}_{\mathrm{dang}}$ is the indicator vector of the dangling nodes, with entries $\mathbf{1}[o(u)=0]$. The third term therefore collects the mass held by nodes without out-edges. With $w(u,v)$ denoting the weight of the directed edge from $u$ to $v$ and $o(u)=\sum_v w(u,v)$ the out-strength, the transition matrix is
\begin{equation}
\label{eq:transition-impl}
P_{u,v}
=
\begin{cases}
w(u,v)/o(u) & \text{if }o(u)>0,\\
0 & \text{otherwise.}
\end{cases}
\end{equation}
The rows of dangling nodes are zero, so $P$ is substochastic; the mass lost at those nodes returns through the third term of Equation~\eqref{eq:pagerank} and is spread according to $\boldsymbol{\pi}$. Filling each zero row with $1/|V|$ gives the row-stochastic matrix of Equation~\eqref{eq:transition}, and with uniform $\boldsymbol{\pi}$ Equation~\eqref{eq:pagerank} is then the Google-matrix form of Equation~\eqref{eq:google-matrix}. The iteration starts from $\mathbf{r}^{(0)}=\mathbf{1}/|V|$ and stops once $\|\mathbf{r}^{(t+1)}-\mathbf{r}^{(t)}\|_1 < \varepsilon$ or after $I_{\max}$ iterations. $V$ is the set of end points of the retained edges, and the final vector is normalized to sum to one over $V$. Contracts can also be listed in descending order of score, ties broken by address, but such ordinal ranks serve only for listing; the correlations of Section~\ref{sec:statistical-tests} are computed on the raw scores, where tied scores stay tied.

The damping factor $d$ interpolates between pure propagation ($d\to 1$) and uniform scores ($d\to 0$). Its default of $0.85$ is the classical web setting \parencite{page1999pagerank}, and the robustness sweep varies it with edge weights held fixed (Section~\ref{sec:scaling-robustness}). The strict $\varepsilon=10^{-8}$ and $I_{\max}=300$ keep the measured runtimes tied to the structure of the graph and not to early termination.

\dissertationSubheading[sec:data-sources]{Data sources and sampling}

This section fixes the chain, the time interval, the public tables, the events and the contracts of the main sample; no later step brings in events from outside it. Figure~\ref{fig:data-overview} places every kind of on-chain record used in the thesis on one calendar, with the months each analysis scores, the months it holds out as labels, and what each record is used for.

\begin{figure}[htbp]
\centering
\begin{adjustbox}{max width=\linewidth}
% !TEX root = ../main.tex
% Blockchain data used in the thesis: one lane per analysis, one bar per kind of
% on-chain record, placed on the months it covers (1 Oct 2023 to 30 Sep 2026).
% The 27 months from October 2023 to December 2025 are drawn at a compressed scale
% (0.1 per month); the nine months of 2026 at 0.8 per month.
\begin{tikzpicture}[x=1cm, y=1cm, font=\scriptsize, oneline/.style={text height=1.6ex, text depth=0.45ex}]
\definecolor{dsApproval}{RGB}{0,114,178}
\definecolor{dsTransfer}{RGB}{0,158,115}
\definecolor{dsGMX}{RGB}{230,159,0}
\definecolor{dsOther}{RGB}{140,140,140}
\def\xa{1.95}% 1 Oct 2023
\def\xj{4.65}% 1 Jan 2026 (start of the uncompressed scale)
\def\xw{7.05}% 1 Apr 2026: W0 freeze t1 (31 Mar 2026, 23:59:59 UTC)
\def\xt{9.45}% 1 Jul 2026: T_obs and W1 freeze (30 Jun 2026, 23:59:59 UTC)
\def\xe{11.85}% 1 Oct 2026: end of the W1 labels
\def\bh{0.36}% bar height
\def\xr{12.45}% left edge of the right-hand column
% \dsInput{colour}{y}{x start}{x end}{text}: records used as input
% (scores, cohorts, same-window checks).
\def\dsInput#1#2#3#4#5{%
  \fill[#1!35] (#3,{#2-\bh/2}) rectangle (#4,{#2+\bh/2});
  \draw[#1!75!black, line width=0.5pt] (#3,{#2-\bh/2}) rectangle (#4,{#2+\bh/2});
  \node[oneline, inner sep=0pt] at ({(#3+#4)/2},#2) {#5};}
% \dsHeld{colour}{y}{x start}{x end}{text}: labels counted after a freeze date
% and kept out of the scores.
\def\dsHeld#1#2#3#4#5{%
  \fill[pattern=north east lines, pattern color=#1] (#3,{#2-\bh/2}) rectangle (#4,{#2+\bh/2});
  \draw[#1!75!black, line width=0.7pt] (#3,{#2-\bh/2}) rectangle (#4,{#2+\bh/2});
  \node[oneline, fill=white, inner sep=1pt] at ({(#3+#4)/2},#2) {#5};}
\def\rowlabel#1#2#3{\node[oneline, anchor=east, text=#1!75!black] at ({\xa-0.1},#2) {#3};}
\def\lanetitle#1#2{\node[oneline, anchor=west, font=\footnotesize\bfseries, fill=white, inner sep=1pt] at (0,#1) {#2};}
\def\laneuse#1#2{\node[anchor=north west, text width=4.9cm, align=flush left, inner sep=0pt] at (\xr,#1) {#2};}
\def\lanerule#1{\draw[black!25] (0,#1) -- (17.4,#1);}
\def\ytop{-0.3}% top of the lanes
\def\ybot{-9.45}% bottom of the lanes

% Calendar: year boundaries on the compressed part, month boundaries in 2026
\foreach \x in {1.95,2.25,3.45} {\draw[black!12] (\x,\ytop) -- (\x,\ybot);}
\foreach \k in {0,...,9} {\draw[black!12] ({\xj+0.8*\k},\ytop) -- ({\xj+0.8*\k},\ybot);}
\foreach \k/\name in {0/Jan,1/Feb,2/Mar,3/Apr,4/May,5/Jun,6/Jul,7/Aug,8/Sep}
  {\node[oneline] at ({\xj+0.8*(\k+0.5)},0) {\name};}
\node[oneline] at ({(\xa+\xj)/2-0.1},0) {compressed scale};
\node[oneline, font=\scriptsize\bfseries] at ({(\xa+\xj)/2},0.34) {Oct 2023 -- Dec 2025};
\node[oneline, font=\scriptsize\bfseries] at ({\xj+0.8*4.5},0.34) {2026};
\draw[black!60] ({\xj-0.09},-0.24) -- ({\xj-0.01},-0.08);
\draw[black!60] ({\xj+0.01},-0.24) -- ({\xj+0.09},-0.08);
\draw[decorate, decoration={brace, amplitude=4pt}] ({\xa+0.03},0.62) -- ({\xt-0.03},0.62)
  node[oneline, midway, above=4pt] {observation window (33 months)};
\draw[decorate, decoration={brace, amplitude=4pt}] ({\xt+0.03},0.62) -- ({\xe-0.03},0.62)
  node[oneline, midway, above=4pt] {W1 label window};

% Freeze dates
\draw[black!70, dashed, line width=0.6pt] (\xw,\ytop) -- (\xw,\ybot);
\draw[black!70, dashed, line width=0.6pt] (\xt,\ytop) -- (\xt,\ybot);
\node[align=center, anchor=north] at (\xw,{\ybot-0.05}) {31 Mar 2026\\W0 freeze $t_1$};
\node[align=center, anchor=north] at (\xt,{\ybot-0.05}) {30 Jun 2026\\$T_{\mathrm{obs}}$, W1 freeze};

% Lane A: cohort selection
\lanetitle{-0.6}{Matched cohort selection}
\rowlabel{dsApproval}{-1.0}{Approval}
\dsInput{dsApproval}{-1.0}{\xa}{\xt}{non-zero approvals from at least three owners}
\laneuse{-0.5}{The $27{,}844$ contracts among the $29{,}022$ spenders that at least three owners approved (\S\ref{sec:data-sources}).}
\lanerule{-1.35}

% Lane B: same-window analysis
\lanetitle{-1.65}{Same-window analysis, matched cohort ($n=27{,}844$)}
\rowlabel{dsApproval}{-2.05}{Approval}
\dsInput{dsApproval}{-2.05}{\xa}{\xt}{EndorseRank; allowance checks}
\rowlabel{dsTransfer}{-2.49}{Transfer}
\dsInput{dsTransfer}{-2.49}{\xa}{\xt}{AWP; transfer and Sybil checks}
\laneuse{-1.55}{Both edge types together give \mbox{C-PR} and \mbox{S-PR}. RQ1, RQ2, the benchmark of RQ3 and the same-window part of RQ5 (\S\S\ref{sec:benchmark-results}--\ref{sec:coupled-results}).}
\lanerule{-2.84}

% Lane C: holdout W0 on spender contracts
\lanetitle{-3.14}{Holdout W0, spender contracts ($n=33{,}101$)}
\rowlabel{dsApproval}{-3.54}{Approval}
\dsInput{dsApproval}{-3.54}{\xa}{\xw}{score input}
\dsHeld{dsApproval}{-3.54}{\xw}{\xt}{new pairs}
\rowlabel{dsTransfer}{-3.98}{Transfer}
\dsInput{dsTransfer}{-3.98}{\xa}{\xw}{score input}
\dsHeld{dsTransfer}{-3.98}{\xw}{\xt}{new senders}
\laneuse{-3.04}{Contract spenders holding an allowance at $t_1$; scores and degree baselines use events up to $t_1$. RQ4 and RQ5 (\S\ref{sec:holdout-protocol}, \S\ref{sec:holdout-results}).}
\lanerule{-4.33}

% Lane D: W0 trader run
\lanetitle{-4.63}{W0 trader run, exploratory ($n=88$)}
\rowlabel{dsApproval}{-5.03}{Approval}
\dsInput{dsApproval}{-5.03}{\xa}{\xw}{score input}
\rowlabel{dsTransfer}{-5.47}{Transfer}
\dsInput{dsTransfer}{-5.47}{\xa}{\xw}{score input}
\rowlabel{dsGMX}{-5.91}{GMX V2}
\dsHeld{dsGMX}{-5.91}{\xw}{\xt}{trader labels}
\laneuse{-4.53}{GMX~V2 contract accounts in the $t_1$ graph with at least three closes from April to June; profit labels under no rule (\S\ref{sub:trader-holdout}). Their liquidation labels enter only the registered comparison.}
\lanerule{-6.26}

% Lane E: registered window W1
\lanetitle{-6.56}{Registered window W1}
\rowlabel{dsApproval}{-6.96}{Approval}
\dsInput{dsApproval}{-6.96}{\xa}{\xt}{score input}
\dsHeld{dsApproval}{-6.96}{\xt}{\xe}{new pairs}
\rowlabel{dsTransfer}{-7.40}{Transfer}
\dsInput{dsTransfer}{-7.40}{\xa}{\xt}{score input}
\dsHeld{dsTransfer}{-7.40}{\xt}{\xe}{new senders}
\rowlabel{dsGMX}{-7.84}{GMX V2}
\dsHeld{dsGMX}{-7.84}{\xt}{\xe}{liquidation-free}
\laneuse{-6.46}{$34{,}587$ spender contracts and $70$ contract traders: rule~B and the comparison on outcomes built from neither edge type. Extracted after the registration of 9~October~2026, 16:49~UTC (\S\ref{sub:fresh-holdout}).}
\lanerule{-8.19}

% Lane F: account type
\lanetitle{-8.49}{Account type}
\rowlabel{dsOther}{-8.95}{Code}
\fill[dsOther!75!black] (12.1,-8.95) +(0,0.13) -- +(0.13,0) -- +(0,-0.13) -- +(-0.13,0) -- cycle;
\node[anchor=east, fill=white, inner sep=1pt] at (11.9,-8.95) {\texttt{eth\_getCode} at pinned blocks, 9 Oct 2026};
\laneuse{-8.39}{Which candidate, spender and trader addresses are contracts; an EIP-7702 delegation designator counts as an EOA (\S\ref{sec:data-sources}).}

% Legend: record colours, then bar styles
\begin{scope}[shift={(0,-10.55)}]
  \node[oneline, anchor=west, font=\scriptsize\bfseries, inner sep=0pt] at (0,0) {Records};
  \foreach \c/\t/\x/\y in {dsApproval/ERC-20 Approval logs/1.4/0, dsTransfer/ERC-20 Transfer logs/5.0/0, dsGMX/GMX V2 position closes/8.6/0, dsOther/contract code read over RPC/12.45/0} {
    \fill[\c!35] (\x,{\y-0.1}) rectangle ({\x+0.45},{\y+0.1});
    \draw[\c!75!black, line width=0.5pt] (\x,{\y-0.1}) rectangle ({\x+0.45},{\y+0.1});
    \node[oneline, anchor=west] at ({\x+0.5},\y) {\t};
  }
  \node[oneline, anchor=west, font=\scriptsize\bfseries, inner sep=0pt] at (0,-0.5) {Bars};
  \fill[dsApproval!35] (1.4,-0.6) rectangle (1.625,-0.4);
  \fill[dsTransfer!35] (1.625,-0.6) rectangle (1.85,-0.4);
  \draw[black!60, line width=0.5pt] (1.4,-0.6) rectangle (1.85,-0.4);
  \node[oneline, anchor=west] at (1.9,-0.5) {solid: input to scores, cohorts or same-window checks};
  \fill[pattern=north east lines, pattern color=black!70] (9.4,-0.6) rectangle (9.85,-0.4);
  \draw[black!60, line width=0.7pt] (9.4,-0.6) rectangle (9.85,-0.4);
  \node[oneline, anchor=west] at (9.9,-0.5) {hatched: held-out labels, counted after the freeze};
  \draw[black!70, dashed, line width=0.6pt] (1.625,-1.1) -- (1.625,-0.74);
  \node[oneline, anchor=west] at (1.9,-0.92) {dashed line: freeze date};
\end{scope}
\end{tikzpicture}

\end{adjustbox}
\caption[{Blockchain data used in this thesis: sources, months covered, held-out labels and uses.}]{Blockchain data used in this thesis, all from Arbitrum One. Each lane is one analysis, and each bar is one kind of on-chain record placed on the months it covers; the 27 months before January 2026 are drawn at a compressed scale. Solid bars are inputs to scores, to cohorts or to same-window checks. Hatched bars are labels counted after a freeze date and kept out of the scores; the July--September logs were extracted only after the registration of W1 was committed at 16:49~UTC on 9~October~2026. The text in a bar says what the records compute, and the right-hand column names the cohort, the research questions and the sections that report the results. The approval logs of five tokens select the matched cohort, and every approval and transfer query is then filtered to it and to those tokens.}
\label{fig:data-overview}
\end{figure}

\begin{itemize}
\item Primary chain: All events come from Arbitrum One (chain ID $42161$), an optimistic rollup on Ethereum \parencite{kalodner2018arbitrum} whose logs are in the public BigQuery export. One chain avoids cross-chain aliases, where the addresses of one entity on several chains would appear as distinct nodes, and every contract in the sample runs under the same fees and latencies.

\item Observation window: Thirty-three months, from 2023-10-01 00:00:00 UTC to 2026-06-30 23:59:59 UTC; the queries take 2026-07-01 as an exclusive end. With the registered label window W1, from July to September 2026, the data cover three years. Both graphs use exactly the same dates. The window is long enough for a contract to collect approvals from several owners and for the tenure and active-month proxies to vary. The time decay is anchored at $T_{\mathrm{obs}}=$ 2026-06-30 23:59:59 UTC, so an event on the final day has $0\le\Delta t<1$ and one on the first day is about $1{,}004$ days old.

\item BigQuery source: Logs are read from the public BigQuery table
\begin{center}
\texttt{bigquery-public-data.goog\_blockchain\_arbitrum\_one\_us.logs}.
\end{center}
One record is one event emitted by a contract, with its block timestamp, block number, log index, emitting address, topics and data. Scan~A reads this table once for the observation window and copies the events the study needs into a table of the study's own Google Cloud project, and scan~B does the same for July to September 2026 (Section~\ref{sec:collection-strategy}). Every later query reads the project tables. Anyone with a Google Cloud account can rerun the same queries against the same public table; Appendix~\ref{ch:appendix-sql} shows the query patterns.

\item Reputation events: EndorseRank uses ERC-20 \texttt{Approval} logs \parencite{vogelsteller2015erc20}, which also cover approvals made through EIP-2612 \texttt{permit} \parencite{lundfall2020permit}, and AWP uses \texttt{Transfer} logs. Both are kept only when the log has exactly three topics, which leaves out ERC-721 tokens (Section~\ref{sec:event-decoding}), and only for the five tokens of Section~\ref{sec:amount-revision}. GMX~V2 position closes supply the trader labels. The plan also listed the borrow, repay and liquidation events of the Aave~V3 pool, and scan~A copied them; no analysis in this thesis uses them.

\item Account type: Whether an address is a contract is read with the \texttt{eth\_getCode} call of a public Arbitrum One RPC endpoint at one pinned block. For the cohort candidates this is block $513{,}216{,}506$, read on 9~October~2026. An address whose code is non-empty is a contract. An address without code is an EOA, and so is an address whose code is an EIP-7702 delegation designator, the prefix \texttt{0xef0100} followed by a 20-byte address: under EIP-7702 an EOA keeps its private key and points to code that runs on its behalf \parencite{buterin2024eip7702}. Holdout spenders and GMX accounts that were not among the cohort candidates were read by the same call at later blocks, within three hours of the first read; each read is stored with its block number.

\item Matched cohort: The cohort $\mathcal{W}$ is every contract that received a non-zero \texttt{Approval} of one of the five tokens from at least three distinct owners inside the observation window. The plan's rule counted approvals of every token: the spender addresses that met it are the candidates, and those with contract code formed the plan's cohort. The revised cohort, $n=|\mathcal{W}|=15{,}107$, is the part of it that meets the rule on the five tokens; Section~\ref{sub:cohorts} gives the account types of all candidates and the size of both cohorts. The plan's cohort was fixed from the approval logs and the code check alone, before any transfer was aggregated, and the revised cohort is fixed from the approvals of the five tokens alone. Because the rule reads approvals, the cohort is not independent of the allowance graph: every cohort contract held allowances at some time in the window, while the rule asks nothing about transfers. A cohort contract need not appear in either graph at $T_{\mathrm{obs}}$. It is outside the allowance graph when every allowance it had received was later set to zero and it held no positive allowance as an owner, and outside the transfer graph when it sent and received no non-zero transfer in the window. The method concerned then scores it $0$ (Section~\ref{sec:graph-construction}); Section~\ref{sub:cohorts} counts these contracts.

\item Ego network: Only events that involve a cohort contract are extracted: an approval in which a cohort contract is the owner or the spender, and a non-zero transfer in which it is the sender or the recipient. Each graph is thus the one-hop ego network of the cohort, since an event between two addresses outside it is never extracted. The plan extracted the events of every token that involve a contract of its cohort; the revised graphs keep those of the five tokens that involve a contract of the revised cohort, a subset, so no log had to be read again. Appendix~\ref{sec:data-events} counts the extracted events, and Section~\ref{sec:dataset} gives the nodes and edges of both graphs at $T_{\mathrm{obs}}$. Owners that are themselves cohort contracts, or that a cohort contract has approved, receive approvals of their own, so the allowance walk can carry endorsement over more than one hop; Section~\ref{sub:theory-closed-form} examines how far it does.

\item Exclusions: Mempool transactions, off-chain order books, centralized exchange accounts and bridges lie outside the data, and no attempt is made to merge addresses controlled by one operator into one entity. Tokens other than the five are left out, and native ether, which moves without \texttt{Transfer} logs, enters only as WETH. No EOA is scored or evaluated.

\end{itemize}

\dissertationSubheading[sec:collection-strategy]{Two-phase collection strategy}

Data collection and evaluation run as two separate phases, so the tables, benchmarks and robustness runs can be rebuilt as often as needed without cloud costs, by us or by anyone who holds the downloaded tables.

\begin{itemize}
\item Phase 1 (BigQuery): Scan~A copies the approval, transfer and GMX events of the observation window from the public logs into a table of the study's project, partitioned by month, and scan~B copies those of July to September 2026 into a second table once the registration of W1 is committed (Section~\ref{sub:fresh-holdout}). One query lists the spenders approved by at least three owners, the code check of Section~\ref{sec:data-sources} fixes the cohort, and a further query extracts the ego events of the cohort. At each anchor date, $T_{\mathrm{obs}}$ and $t_1$, the ego events are then aggregated into pair tables. For the allowance graph the latest allowance of each (token, owner, spender) triple is taken in (block number, log index) order, kept when it is positive, and summed into one row per owner--spender pair that carries the EndorseRank weight and the capped US-dollar amount. For the transfer graph each sender--recipient pair gets one row with its AWP weight, its US-dollar weight, its largest time weight and its number of transfers. In the revised analysis these queries run on the ego events of the five tokens, valued with a table of daily prices loaded into the project. The same phase computes the same-window proxies and the holdout labels and maps every address to an integer node identifier.

\item Phase 2 (local evaluation): The pair tables, the node table, the proxies, the labels and the GMX closes are downloaded once. From them a local machine decodes the GMX closes and builds the graphs, and it computes the scores, the checks, the alignment statistics, the holdout evaluations, the benchmark and the robustness runs. No step of Phase~2 queries BigQuery.

\item Cost guardrails: Every query is first run as a dry run, and its estimated cost is checked against the budget of USD~50 that the configuration sets for the whole study. A query whose estimate would take the spending past the budget is refused unless the operator confirms it with an extra flag.\footnote{The budget is enforced by the study's own pipeline; BigQuery does not impose it.} Every job is written to a cost ledger with the bytes it billed and its cost, so the cost of acquisition can be audited (Appendix~\ref{ch:appendix-eval}).

\item Storage: Every file written under the data folder of the repository is split into parts below 90~MB, so that the extracts and summaries can be kept in Git under GitHub's file-size limit.

\end{itemize}

\dissertationSubheading[sec:pipeline-architecture]{Pipeline architecture}

Figure~\ref{fig:pipeline-architecture} traces the whole pipeline, from public logs to the research tables. All remote steps write tables in the study's BigQuery project, and the downloaded parts are not changed afterwards. Given its inputs and the versioned configuration file, every local step is deterministic.

\begin{figure}[htbp]
\centering
\begin{adjustbox}{max width=\linewidth}
\begin{tikzpicture}[
  font=\small,
  box/.style={draw, rounded corners=2pt, align=center, minimum height=0.95cm, minimum width=2.2cm, inner sep=3pt},
  wide/.style={draw, rounded corners=2pt, align=center, minimum height=0.95cm, minimum width=2.6cm, inner sep=3pt},
  arrow/.style={-{Stealth[length=2.5mm]}, thick}
]
% Top row: remote acquisition (BigQuery and one RPC call), left to right
\node[wide] (bq) at (0,0) {BigQuery\\public logs};
\node[wide] (scan) at (3.6,0) {Scan A, scan B\\(copy approvals,\\transfers, GMX closes)};
\node[wide] (cohort) at (7.3,0) {Cohort rule\\($\geq 3$ owners, five tokens;\\\texttt{eth\_getCode})};
\node[wide] (pairs) at (11.0,0) {Ego events in US dollars;\\pair tables at $T_{\mathrm{obs}}$, $t_1$;\\proxies, labels};

% Local evaluation. Row 1: parts -> graphs -> scores -> evaluation -> export.
% Rows 2 and 3: proxies and labels, and the decoded GMX closes, feed the evaluation.
\node[box] (parts) at (0,-2.5) {Parquet parts\\(below 90\,MB)};
\node[box] (graphs) at (3.1,-2.5) {Build\\graphs};
\node[wide] (ranks) at (6.4,-2.5) {PageRank scores\\(EndorseRank, AWP,\\C-PR, S-PR)};
\node[wide] (eval) at (10.3,-2.5) {Same-window checks,\\holdouts W0 and W1,\\benchmark};
\node[wide] (tabexp) at (13.9,-2.5) {Tables \&\\figures export};
\node[box] (proxy) at (3.1,-4.1) {Proxies and\\holdout labels};
\node[box] (gmx) at (3.1,-5.5) {Decode GMX closes;\\trader labels};

\draw[arrow] (bq) -- (scan);
\draw[arrow] (scan) -- (cohort);
\draw[arrow] (cohort) -- (pairs);
\draw[arrow] (pairs.south) -- ++(0,-0.5) -| (parts.north);
\draw[arrow] (parts) -- (graphs);
\draw[arrow] (graphs) -- (ranks);
\draw[arrow] (ranks) -- (eval);
\draw[arrow] (eval) -- (tabexp);
\draw[arrow] ([xshift=4pt]parts.south) |- (proxy.west);
\draw[arrow] ([xshift=-4pt]parts.south) |- (gmx.west);
\draw[arrow] (proxy.east) -| ([xshift=-8pt]eval.south);
\draw[arrow] (gmx.east) -| ([xshift=8pt]eval.south);

\node[draw, dashed, rounded corners, fit=(bq)(scan)(cohort)(pairs), inner sep=8pt, label=above:{\footnotesize Remote: BigQuery (dry run, budget, cost ledger) and an Arbitrum One RPC endpoint}] {};
\node[draw, dashed, rounded corners, fit=(parts)(graphs)(ranks)(eval)(tabexp)(proxy)(gmx), inner sep=8pt, label=below:{\footnotesize Local / reproducible}] {};
\end{tikzpicture}
\end{adjustbox}
\caption[{Pipeline overview: the study's events are copied and aggregated in BigQuery, the cohort is fixed, the pair tables are downloaded, and the scores, checks, holdouts and benchmark are computed locally.}]{Pipeline overview. In BigQuery, scan~A (observation window) and scan~B (W1, run only after its registration was committed) copy the study's events from the public logs into a project table; the cohort is fixed from the approvals of five tokens and the account code; the ego events of the cohort are valued in US dollars with daily token prices and aggregated into pair tables at each anchor date, together with the proxies and the holdout labels. Every query is dry-run against the study's budget and written to a cost ledger. Locally, the downloaded parts become the graphs from which EndorseRank, AWP, C-PR and S-PR are computed, and the GMX closes are decoded into trader labels. The same-window checks, the holdouts and the benchmark follow, and the tables and figures are exported last.}
\label{fig:pipeline-architecture}
\end{figure}

The Phase~1 queries write the pair tables, the proxies and the labels, which the Phase~2 drivers read to build the graphs, compute the scores and write the machine-readable summaries behind the tables of Chapter~\ref{ch:results}. Unit tests run the scoring and statistics code on small synthetic fixtures.

\dissertationSubheading[sec:event-decoding]{Event decoding}

A log emitted by a smart contract has topics and a data payload but no field names. Decoding turns each row into a typed record for graph construction and the checks, and an error here would corrupt both graphs silently. The topic hashes and field mappings are listed below.

\begin{itemize}
\item Log anatomy: Each log records the emitting contract, a sequence of topics and an ABI-encoded data blob. The first topic is the keccak-256 hash of the event signature; indexed parameters follow as $32$-byte words (addresses left-padded), and non-indexed parameters go into the data field in ABI order \parencite{vogelsteller2015erc20}. ERC-721 tokens emit \texttt{Transfer} and \texttt{Approval} under the same signatures but index the token ID as a fourth topic, so the condition of exactly three topics keeps them out. The ERC-20 events are decoded in SQL during scan~A and the ego extraction; the GMX events are decoded locally.

\item ERC-20 \texttt{Approval}: Its topic hash is
\begin{center}
\texttt{0x8c5be1e5ebec7d5bd14f71427d1e84f3dd0314c0f7b2291e5b200ac8c7c3b925},
\end{center}
the keccak-256 hash of \texttt{Approval(address,address,uint256)}. Owner and spender occupy Topics~$1$ and~$2$, and the allowance is decoded from the data field in the token's base units and valued in US dollars (Section~\ref{sec:amount-revision}). An unlimited allowance (\texttt{type(uint256).max}) keeps its value, far above any other amount, and an allowance of zero is kept, since it sets the allowance in force to zero. A \texttt{permit} emits the same \texttt{Approval} event, so an allowance counts whether it came from an on-chain \texttt{approve()} call or from an off-chain signature \parencite{lundfall2020permit}.

\item ERC-20 \texttt{Transfer}: Its topic hash is
\begin{center}
\texttt{0xddf252ad1be2c89b69c2b068fc378daa952ba7f163c4a11628f55a4df523b3ef},
\end{center}
the keccak-256 hash of \texttt{Transfer(address,address,uint256)}. Sender and recipient are indexed, and the value is decoded in base units and valued in US dollars at the price of its day. Zero-value transfers are dropped, and so are transfers made before their token's first price, which are counted (Appendix~\ref{sec:data-tokens}). Mint and burn logs, which use the zero address, are kept whenever a cohort contract is involved. Self-transfers stay in the events but are left out of the inbound-counterparty computations (Section~\ref{sec:proxy-formulas}), since they reveal no counterparty.

\item GMX~V2 position closes: A close is one \texttt{EventLog1} log of the GMX~V2 event emitter whose event name is \texttt{PositionDecrease}, so partial closes and liquidations count. The trading account is an indexed topic, so the accounts can be listed in SQL; the data field is decoded locally with the emitter's ABI into the account, the base profit or loss in USD, the order type and the size of the close. A close is a liquidation when its order type is the liquidation type or when the liquidation handler emitted it. The decoder read every close of a contract account in the observation window and in W1 without a decoding error (Section~\ref{sub:cohorts} gives the counts).

\item Address normalization and units: Addresses are compared as 20-byte values and written in lower-case hexadecimal. Amounts are the \texttt{uint256} value as a double-precision float, divided by ten to the power of the token's decimals and multiplied by its US-dollar price (Section~\ref{sec:amount-revision}). An unlimited allowance, $2^{256}-1\approx1.16\times10^{77}$ base units, stays far above any other amount; Sections~\ref{sub:endorserank-graph} and~\ref{sub:coupled-graph} give how each score treats it. A data field that cannot be read as a number gives the amount zero, so such a transfer is dropped and such an approval counts as setting the allowance to zero.

\end{itemize}

\dissertationSubsubheading[sec:amount-revision]{Tokens, prices and the revision of the amounts}

The analysis plan fixed amounts in raw base units, every ERC-20 token, and the value transform of Equation~\eqref{eq:value-transform} with $b=1$. Three defects of that setting came to light after the registered results of the same window, W0 and W1 had been computed. First, $V(z)$ is within $10^{-4}$ of $1$ for any amount of $10$ base units or more, and almost every transfer and latest allowance reaches that level (Appendix~\ref{sec:data-tokens}), so the amount played no part in EndorseRank or AWP, although AWP's value weight exists so that larger transactions count for more \parencite{do2023awp}. Second, the layers of C-PR added raw base units across tokens: one USDC is $10^{6}$ base units and one WETH $10^{18}$, and an unlimited allowance, $2^{256}-1$, took almost the whole row of every owner that held one. Third, every ERC-20 contract counted. A token contract can emit \texttt{Approval} and \texttt{Transfer} logs that name any address as owner or sender without that address taking part, as spam and address-poisoning tokens do, so a log of such a token is no evidence that the named owner endorsed or paid anyone.

The amounts were therefore revised on 10~October~2026, and the revised analysis is the one that Chapter~\ref{ch:results} reports. The revision is a deviation from the plan and from the registration of W1, decided after their results were known, so none of its results is a registered result. It was written down and committed to the public repository before any revised score, proxy or label was computed: the revision document at 07:54~UTC (commit \texttt{a897475}), the choice of five tokens at 08:12~UTC (commit \texttt{09dfc32}) and an amendment of the allowance cap described below at 11:24~UTC (commit \texttt{2705b28}). The registered analysis is kept as it was obtained; its decision contrasts and verdicts are set beside the revised ones in Appendix~\ref{sec:map-registered}, and its complete results are published unchanged (Appendix~\ref{ch:appendix-eval}).

\begin{itemize}
\item Token list: A token qualified when Chainlink's directory of data feeds on Arbitrum One, read on 10~October~2026, has a market-price feed for it, quoted in US dollars or in ether, whose market-risk category is low or medium, and when that feed prices the token itself, the asset that the token wraps one-to-one (WETH for ether), or the Ethereum token of which the token is the copy minted by Arbitrum's canonical bridge. Chainlink assigns the category from the depth, spread and concentration of an asset's markets, and leaves exchange-rate and other custom feeds, which do not follow the market, outside it \parencite{chainlinknddatafeeds}. Of the 133 Arbitrum One tokens that the directory maps, 64 qualified and 46 of those have a daily price. The analysis keeps the five of these with the most rows in the data of the observation window: WETH, USDC, USDT0, ARB and WBTC. All five are priced by feeds rated low risk. Appendix~\ref{sec:data-tokens} gives their addresses and the rows that every other token holds.
\item Prices: The US-dollar price of a token on a UTC day is its closing price in the daily series of the DefiLlama coins API \parencite{defillamanddocs}, the point nearest to 00:00~UTC of the following day. Where the series under the token's Arbitrum address starts later than the token's use (for WETH, in August 2024), the days are filled from DefiLlama's series of the same asset under its CoinGecko identifier, provided the two series differ by less than $1\%$ in median on the days they share. A daily point more than twice or less than half of both neighbouring days is treated as a data error and replaced by the previous day's price; a day without a point takes the previous day's price. For a sample of dates, the prices are compared with the Chainlink feeds read on the chain (Appendix~\ref{sec:data-tokens}).
\item Amounts: A transfer of $v$ base units of a token with $q$ decimals, on a day with price $p$, is worth $x_e=v\,10^{-q}p$ US dollars. A latest allowance is valued at the price of the anchor's day, 30~June~2026 at $T_{\mathrm{obs}}$ and 31~March~2026 at $t_1$, since the edge stands for the permission in force at the anchor. Native ether moves without \texttt{Transfer} logs and was never in the data; it enters the graphs through WETH only.
\item Value slopes: The slope of $V$ is set so that the median amount weighs one half, $b=\ln 3/m$, with $m$ the median amount rounded to two significant digits. For transfers, $m_T$ is the median US-dollar value of the ego transfers up to $t_1$, $m_T=200$ dollars; for allowances, $m_A$ is the median value at $t_1$ prices of the finite positive latest allowances at $t_1$, $m_A=62$ dollars. Both use events up to the earlier freeze only, so no label window enters them.
\item Allowance cap of C-PR: In the allowance layer of C-PR an allowance counts at most $U=20{,}000$ dollars, the 99th percentile of the US-dollar values of single transfers up to $t_1$, rounded to two significant digits. The revision document first defined $U$ as the 99th percentile of the finite latest allowances at $t_1$. On the five tokens that percentile is about $10^{47}$ dollars, because the allowances below $2^{256}-1$ include many that are unlimited in effect, such as $2^{255}$ or $2^{128}$, and a cap at that level would have left unlimited allowances in charge of every owner's row. The definition was amended before any revised score existed; the medians were unchanged.
\item Scope: Every quantity of the revised analysis uses the five tokens only. The matched cohort is the set of contracts that received a non-zero \texttt{Approval} of one of the five tokens from at least three distinct owners in the observation window. Such a contract was approved by three owners on some token, so it belongs to the cohort of the plan and its account type was already known; the graphs, the proxies, the spender cohorts and the spender labels are drawn from the ego events of the plan's cohort, filtered to the five tokens and to the revised cohort, and no public log was scanned again. The GMX trader labels do not depend on tokens.
\end{itemize}

On the revised scores, rules A and B and the comparison on outcomes built from neither edge type are applied with the thresholds and the wording that were fixed for them, $\delta=0.01$ included (Sections~\ref{sec:holdout-protocol}, \ref{sub:fresh-holdout} and~\ref{sub:neutral-label-protocol}). Because the scores they judge were defined after the registered results were known, these verdicts are a re-analysis and not registered tests; Chapter~\ref{ch:results} says so wherever it reports them.

\dissertationSubheading[sec:graph-construction]{Graph construction}

EndorseRank and AWP weight their edges the same way and run the same PageRank solver, Equation~\eqref{eq:pagerank}; they differ in the events that make the edges and in where the walk restarts. Both graphs hold the ego network of the matched cohort $\mathcal{W}$ (Section~\ref{sec:data-sources}), which brings in the cohort's one-hop neighbors as nodes. A method reports a score for every contract in the evaluated set $\mathcal{C}$: the matched cohort in the same window, and the spender and trader cohorts in the holdouts. A contract missing from the graph scores $0$.

Edges in EndorseRank stand for endorsements and edges in AWP for transfers (Section~\ref{sec:allowance-endorsement}). Both methods weight an event by its age, so that recent events count for more, and pass its US-dollar amount through a bounded function, so that a larger amount counts for more but one event counts at most once. AWP is described first, since EndorseRank takes its weights from it.

\dissertationSubsubheading[sub:awp-graph]{AWP transfer graph}

AWP was proposed by \textcite{do2023awp} at IEEE RIVF 2023. A transfer $e$ worth $x_e$ US dollars, made $\Delta t_e$ days before the anchor ($T_{\mathrm{obs}}$, or $t_1$ in W0), enters the edge from its sender to its receiver with the weight $\sigma(\Delta t_e)\,V(x_e)$. The time weight is the logistic decay
\begin{equation}
\label{eq:logistic-decay}
\sigma(\Delta t)
=
\frac{1}{1+\exp\bigl(k(\Delta t-t_0)\bigr)},
\qquad
k=0.01,\quad t_0=180\text{ days},
\end{equation}
and the value weight is
\begin{equation}
\label{eq:value-transform}
V(z)
=
\frac{2}{1+e^{-bz}}-1,
\qquad
b=\frac{\ln 3}{m},
\end{equation}
which is $0$ at $z=0$, $1/2$ at $z=m$ and rises towards $1$ as the amount grows. For transfers $m=m_T=200$ dollars and for allowances $m=m_A=62$ dollars, the medians of Section~\ref{sec:amount-revision}, so $b_T=\ln 3/200$ and $b_A=\ln 3/62$ per dollar. The edge weights are sums over transfers:
\begin{equation}
\label{eq:transfer-weight}
w_T(u,v)
=
\sum_{e:\,u\to v}
\sigma(\Delta t_e)\,V(x_e).
\end{equation}
The sum runs over every transfer of the five tokens from $u$ to $v$ up to the anchor. A self-transfer ($u=v$) is kept and becomes a self-loop on $u$. The walk restarts at an address in proportion to its activeness,
\begin{equation}
\label{eq:activeness}
X_u
=
\sum_{v\neq u}\;
\max_{e:\,u\to v}\sigma(\Delta t_e),
\end{equation}
which adds up, over the addresses that $u$ has sent to, the time weight of its latest transfer to each. The restart distribution of Equation~\eqref{eq:pagerank} is $\pi_u=X_u/\sum_{v\in V_T}X_v$, and it also takes the mass of the dangling nodes. An address that never sends to another has $X_u=0$ and receives no restart mass, so its score comes from its in-edges alone. The damping factor is $d=0.85$.

The paper scores Ethereum transactions and gives no parameter values; its time weight is the logistic function of Equation~\eqref{eq:logistic-decay} written in calendar time instead of age. The values of $k$ and $t_0$ are this study's choices, fixed in the plan; the plan fixed $b=1$ on raw base units, and the revision replaced it (Section~\ref{sec:amount-revision}). With $b_T=\ln 3/m_T$ a transfer of the median amount weighs one half, one of twice the median $0.8$ and one of half the median $0.27$, so an edge weight grows with the amounts transferred as well as with their number. The activeness of an address still counts the addresses it has sent to, each by the time weight of its latest transfer.

Over the 33 months of the window the decay discounts old events strongly (Figure~\ref{fig:logistic-decay-methods}). The weight falls from $\sigma(0)=0.858$ through $\sigma(90)=0.711$, $\sigma(182)=0.495$ and $\sigma(365)=0.136$ to $\sigma(730)=0.004$, and an event from the first day of the window, $1{,}004$ days before $T_{\mathrm{obs}}$, weighs $\sigma(1{,}004)=0.0003$. Events older than about a year still create edges but add little weight. Because each row of $P$ is divided by the out-strength, the decay changes the walk only where one address has out-edges of different ages: an owner whose approvals are all old passes on as much mass as one whose approvals are recent, and only the split of that mass among its edges follows their ages. In AWP the decay also enters the restart weights, so senders that were active only long ago receive little restart mass.

\begin{figure}[htbp]
\centering
\includegraphics[width=0.8\textwidth]{\figpath{logistic-decay.pdf}}
\caption[{Logistic time-decay $\sigma(\Delta t)$ of AWP and EndorseRank ($k=0.01$, $t_0=180$ days). Horizontal axis is the age of an event in days before the anchor; vertical axis is the time weight of a transfer or approval.}]{Logistic time-decay $\sigma(\Delta t)$ of AWP and EndorseRank ($k=0.01$, $t_0=180$ days). The horizontal axis shows the age of an event in days before the anchor, $T_{\mathrm{obs}}=$ 30~June~2026 or, in W0, $t_1=$ 31~March~2026, and the vertical axis shows the time weight of a transfer or an approval. The shaded band marks the ages that occur at $T_{\mathrm{obs}}$ (0 to $1{,}004$ days) and the arrow those at the W0 freeze (0 to 913 days). Over the band the weight falls from $0.86$ to below $0.001$.}
\label{fig:logistic-decay-methods}
\end{figure}

Algorithm~\ref{alg:awp} states the construction and scoring steps formally.

\begin{algorithm}[htbp]
\caption{AWP graph construction and scoring}
\label{alg:awp}
\begin{algorithmic}[1]
\Require Transfers $\mathcal{T}$ with US-dollar amounts, seed contracts $\mathcal{W}$, evaluated contracts $\mathcal{C}$, anchor $T$, parameters $k,t_0,b_T,d,\varepsilon,I_{\max}$
\Ensure Scores $\mathbf{s}_{\mathrm{AWP}}$ on $\mathcal{C}$
\State $E \gets \emptyset$; $X_u\gets 0$ for every address $u$
\For{each transfer $e=(u,v,x_e,t_e)$ in $\mathcal{T}$ with $t_e\le T$}
  \State $\Delta t \gets (T-t_e)$ in days
  \State accumulate $\sigma(\Delta t)\,V(x_e)$ on edge $(u\to v)$ \Comment{Equations~\eqref{eq:logistic-decay}--\eqref{eq:transfer-weight}; $u=v$ gives a self-loop}
\EndFor
\For{each sender $u$}
  \State $X_u \gets \sum_{v\neq u}\max_{e:\,u\to v}\sigma(\Delta t_e)$ \Comment{Equation~\eqref{eq:activeness}}
\EndFor
\State remove non-positive edges; retain edges incident to $\mathcal{W}$
\State $\pi_u \gets X_u/\sum_{v}X_v$ for the nodes $u$ of the retained edges
\State $\mathbf{s} \gets \mathrm{WeightedPageRank}(E, \boldsymbol{\pi}, d, \varepsilon, I_{\max})$ \Comment{Equations~\eqref{eq:pagerank}--\eqref{eq:transition-impl}}
\State \Return $\mathbf{s}_{\mathrm{AWP}}[i] \gets \mathbf{s}[i]$ for $i\in\mathcal{C}$ (else $0$)
\end{algorithmic}
\end{algorithm}

\dissertationSubsubheading[sub:endorserank-graph]{EndorseRank allowance graph}

EndorseRank puts AWP's edge weights on the allowance graph. Let $\mathcal{A}$ denote the set of latest allowances after preprocessing. For each (token, owner, spender) triple only the most recent \texttt{Approval} event up to the anchor is kept, in (block number, log index) order, and a zero allowance deletes the edge. The surviving allowance $a^{\star}(u,v,\mathrm{token})$, valued in US dollars at the anchor's price, was set $\Delta t^{\star}(u,v,\mathrm{token})$ days before the anchor, and it enters the edge as one event of that age and amount:
\begin{equation}
\label{eq:approve-weight}
w_A(u,v)
=
\sum_{\mathrm{token}}
\sigma\bigl(\Delta t^{\star}(u,v,\mathrm{token})\bigr)\,
V\bigl(a^{\star}(u,v,\mathrm{token})\bigr),
\end{equation}
with $\sigma$, $V$, $k$ and $t_0$ as for AWP and the slope $b_A=\ln 3/m_A$. A directed edge $u\to v$ means that owner $u$ endorses spender $v$. An allowance of the median amount enters with half its time weight and an unlimited allowance with its full time weight, so $w_A(u,v)$ grows with the amounts that $u$ has approved to $v$, but no single allowance counts for more than one. The damping factor is $d=0.85$.

The walk restarts uniformly. AWP's restart rule, carried over to allowances, would restart the walk at owners in proportion to the number of spenders they approved recently. That rewards granting endorsements, while EndorseRank is meant to rank the contracts that receive them. The analysis plan fixed uniform restarts before any data were extracted and named the variant with AWP's restarts as one to report beside EndorseRank (Table~\ref{tab:endorserank-restarts}).

The weight keeps the amount and the time of the latest approval; it is not the allowance in force at the anchor. Within a pair, on each token, the newest approval replaces the older ones \parencite{vogelsteller2015erc20}. Approvals granted before 1~October~2023 are not observed, a finite allowance that \texttt{transferFrom} has since reduced keeps its approved value, and an owner who approves an allowance-manager contract such as Uniswap's Permit2 appears only with an edge to that contract, not to the application that later spends through it.

The steps for building the graph and running PageRank are listed in Algorithm~\ref{alg:endorserank}.

\begin{algorithm}[htbp]
\caption{EndorseRank graph construction and scoring}
\label{alg:endorserank}
\begin{algorithmic}[1]
\Require Latest allowances $\mathcal{A}$ up to the anchor in US dollars, seed contracts $\mathcal{W}$, evaluated contracts $\mathcal{C}$, anchor $T$, parameters $k,t_0,b_A,d,\varepsilon,I_{\max}$
\Ensure Scores $\mathbf{s}_{\mathrm{ER}}$ on $\mathcal{C}$
\State $E \gets \emptyset$
\For{each latest allowance $(u,v,\mathrm{token},a^{\star},t^{\star})$ in $\mathcal{A}$}
  \State $\Delta t^{\star} \gets (T-t^{\star})$ in days
  \State accumulate $\sigma(\Delta t^{\star})\,V(a^{\star})$ on edge $(u\to v)$ \Comment{Equation~\eqref{eq:approve-weight}}
\EndFor
\State remove non-positive edges; retain edges incident to $\mathcal{W}$
\State $\mathbf{s} \gets \mathrm{WeightedPageRank}(E, \text{uniform }\boldsymbol{\pi}, d, \varepsilon, I_{\max})$ \Comment{Equations~\eqref{eq:pagerank}--\eqref{eq:transition-impl}}
\State \Return $\mathbf{s}_{\mathrm{ER}}[i] \gets \mathbf{s}[i]$ for $i\in\mathcal{C}$ (else $0$)
\end{algorithmic}
\end{algorithm}

Approvals are rarer than transfers, so at $T_{\mathrm{obs}}$ the allowance graph of the cohort has fewer edges than its transfer graph (Section~\ref{sec:dataset} gives both counts). Since one PageRank iteration costs $O(|E|)$ \parencite{leskovec2020mining}, this sparsity is the structural basis of H3, and any difference in runtime that follows from it is a property of the graph and not of the algorithm.

For both methods, edge lists are stored as (from, to, weight) triples; duplicate edges are merged by summing their weights, and non-positive weights are dropped before the solve. The solver builds the transposed transition matrix in compressed sparse-row form and the restart vector with vectorised array operations, iterates Equation~\eqref{eq:pagerank} and returns a score for each node. The scores of the evaluated contracts are read from that vector, with zeros for contracts outside the graph.

\dissertationSubsubheading[sub:coupled-graph]{Coupled operator over both layers (C-PR) and the seeded ablation (S-PR)}

EndorseRank and AWP each walk on a single edge type, and RQ5 asks whether one walk over both edge types ranks contracts better than a walk over one. The analysis plan fixed the coupled operator on 9~October~2026, before any log was extracted, on two layers that weight each event by its amount; the revision values the amounts in US dollars and caps the allowances (Section~\ref{sec:amount-revision}):
\begin{equation}
\label{eq:raw-layers}
\tilde w_A(u,v)
=
\sum_{\mathrm{token}}
\min\bigl(a^{\star}(u,v,\mathrm{token}),U\bigr),
\qquad
\tilde w_T(u,v)
=
\sum_{e:\,u\to v}
x_e\,\sigma(\Delta t_e).
\end{equation}
The allowance layer has the edges of EndorseRank and the transfer layer those of AWP; only the weights differ. Amounts enter in US dollars without the bounded transform, an allowance counts at most $U=20{,}000$ dollars (Section~\ref{sec:amount-revision}), and the allowance layer carries no time decay. Both decision rules for RQ5 were fixed in the plan on the same layers in raw base units; here they are applied to the layers in US dollars. The two layers form one graph on the union node set $V=V_A\cup V_T$, with transition matrices $P_A$ and $P_T$ from Equation~\eqref{eq:transition-impl}. For node $u$ the coupled row is a convex mixture of the two layer rows, renormalized over the layers in which $u$ has out-edges:
\begin{align}
P^{(\lambda)}_{u,\cdot}
&=
\alpha_A(u)\,P_{A,u,\cdot}+\alpha_T(u)\,P_{T,u,\cdot},
\label{eq:coupled-transition}\\
\alpha_A(u)
&=
\frac{\lambda\,\mathbf{1}[o_A(u)>0]}{D_{\lambda}(u)},
\qquad
\alpha_T(u)=\frac{(1-\lambda)\,\mathbf{1}[o_T(u)>0]}{D_{\lambda}(u)},
\nonumber
\end{align}
where $D_{\lambda}(u)=\lambda\,\mathbf{1}[o_A(u)>0]+(1-\lambda)\,\mathbf{1}[o_T(u)>0]$, and $o_A(u)$ and $o_T(u)$ are the out-strengths of $u$ in the two layers. A node with out-edges in only one layer follows that layer with weight one. A layer of weight $0$ is dropped with its edges and nodes, and a node with $D_{\lambda}(u)=0$ is dangling, its mass handled by the third term of Equation~\eqref{eq:pagerank}. Teleportation is uniform, $d=0.85$, and $\mathbf{s}_{\mathrm{CPR}}$ is the PageRank vector of Equation~\eqref{eq:pagerank} with $P=P^{(\lambda)}$, under the solver settings of the single-layer methods.

At $\lambda=1$ the operator walks the allowance layer alone and at $\lambda=0$ the transfer layer alone; we write these walks C-PR ($\lambda=1$) and C-PR ($\lambda=0$). They are the comparators of both rules for RQ5 (Sections~\ref{sec:holdout-protocol} and~\ref{sub:fresh-holdout}). Unit tests check on small fixtures that each end point gives the scores of its layer walked alone with uniform restarts. The family is not continuous at the end points. For $0<\lambda<1$ the node set is $V_A\cup V_T$ and a node with out-edges in only one layer follows it whatever $\lambda$ is, whereas at $\lambda=1$ a node whose out-edges are all transfers is dangling or absent; as $\lambda$ approaches $1$, C-PR therefore tends to a walk that still follows the transfer edges of such nodes, and the same holds at the other end. The values $\lambda=0.5$ (primary) and $\lambda\in\{0.25,0.75\}$ (sensitivity) were fixed in the configuration of 9~October~2026.

In raw base units the plan's layers were very uneven: an owner that held an unlimited allowance passed almost all of its mass through it, and in the transfer layer a token with many decimals outweighed one with few (Section~\ref{sec:amount-revision}). In US dollars the transfer layer adds comparable amounts, and the cap $U$ lets an unlimited allowance weigh as much as an allowance of the size that 99\% of transfers do not exceed, so an owner that holds unlimited and finite allowances side by side splits its row among both (Appendix~\ref{sec:data-tokens} gives the counts). EndorseRank and AWP bound every amount through $V$ instead.

In the S-PR ablation the transfer layer is walked, but the uniform teleportation vector $\boldsymbol{\pi}$ of Equation~\eqref{eq:pagerank} is replaced by the score of C-PR ($\lambda=1$), restricted to $V_T$ and renormalized over it. Let $\mathbf{r}_{A}$ be the PageRank vector of the allowance layer walked alone on $V_A$, extended by $r_{A,u}=0$ for $u\notin V_A$. Then
\[
\pi_u=\frac{r_{A,u}}{\sum_{v\in V_T} r_{A,v}}\qquad(u\in V_T),
\]
so a transfer-graph node without an allowance score receives zero restart mass. If no node of $V_T$ lies in $V_A$, the denominator is zero and $\boldsymbol{\pi}$ falls back to the uniform distribution on $V_T$. The same vector redistributes dangling mass. The result is the TrustRank construction \parencite{gyongyi2004trustrank}, but with allowance-derived seeds instead of hand-labeled ones. Allowances decide where the walk restarts; transfers decide how it spreads. A comparison of S-PR with C-PR separates what the allowance edge contributes as a propagation medium from what it contributes as a seed.

The edge set determines what the coupled operator costs. Because C-PR walks on $E_A\cup E_T$, one of its iterations costs as much as an iteration on each layer, and its runtime is expected to be at or above AWP's. S-PR needs one solve on each layer. Claim~C1 therefore concerns the allowance graph that EndorseRank walks; the coupled walks carry the edge count of the transfer layer with them. Both operators are timed under the benchmark protocol of Section~\ref{sec:computational-benchmark}.

\dissertationSubheading[sec:holdout-protocol]{Temporal holdout protocol}

In the holdout W0, scores use every extracted event from the start of the data, 1~October~2023, up to $t_1=$ 31~March~2026, 23:59:59~UTC. Labels are counted on events from 1~April~2026, 00:00:00~UTC, up to $t_2=$ 30~June~2026, 23:59:59~UTC. The spender cohort consists of the contracts that hold at least one positive latest allowance of one of the five tokens at $t_1$ in the extracted approval data; spenders without code or with an EIP-7702 designator are left out (Section~\ref{sub:cohorts} gives the counts). It includes spenders outside the matched cohort, which entered the extract because a cohort contract approved them. Every score (EndorseRank with both restart rules, AWP, C-PR at three values of $\lambda$ and at its two end points, S-PR and two raw-degree baselines) is computed from events up to $t_1$ alone. The time decay is anchored at $t_1$ as well: $\Delta t$ in Equation~\eqref{eq:logistic-decay} is measured back from $t_1$, so ages run up to 913 days and the weights from $0.858$ down to below $0.001$.

We count two labels per spender, one for each construct, so that in the future window each single-layer method is tested on its own construct and the coupled walks on both:
\begin{enumerate}
\item New approval pairs: the number of owners that give the spender a positive approval in the label window although their owner--spender pair held no positive allowance at $t_1$. This label measures the endorsement construct.
\item New transfer senders: the number of distinct addresses that send the spender a transfer in the label window without having sent it one between the data start (1~October~2023) and $t_1$, self-transfers excluded. It measures the flow construct and is the transfer-side counterpart of the first label.
\end{enumerate}
Both labels count approvals and transfers of the five tokens only.
For baselines we take each spender's in-approve degree (the number of distinct owners with a positive latest allowance) and its transfer in-degree (the number of distinct senders other than the spender itself), both at $t_1$. Each single-layer PageRank smooths one of these local statistics. Because their rows appear in the table, the increment of each PageRank over its own raw degree can be read off directly. That increment answers RQ4; the PageRank coefficient on its own does not.

Uncertainty comes from the paired bootstrap over spender contracts of Section~\ref{sec:statistical-tests} ($B=400$, one shared index matrix), so every difference between two coefficients on the same label is a paired difference. The two contrasts fixed in the plan before any log was extracted are those of rule~A below, applied here to the revised scores after the registered results were known: C-PR minus C-PR ($\lambda=1$) on new approval pairs, and C-PR minus C-PR ($\lambda=0$) on new transfer senders. Every other contrast in this window was computed after the labels were known: each single-layer PageRank against its own degree, C-PR ($\lambda=1$) against the in-approve degree, the two contrasts of S-PR with the layers walked alone, C-PR against C-PR ($\lambda=0$) on new approval pairs, and C-PR against EndorseRank and AWP. The differences between EndorseRank and AWP are in Appendix~\ref{ch:appendix-er-awp}. AWP scores zero a spender without a transfer edge at $t_1$, so the degree contrasts are also computed on the spenders that had received a transfer before $t_1$ (Table~\ref{tab:holdout-receiving}); this restriction was not fixed in advance either.

Rule~A, fixed in the plan, judges C-PR against its own layers walked alone and has two parts. The primary part uses W0: C-PR at $\lambda=0.5$ should do no worse than C-PR ($\lambda=1$) on new approval pairs (the interval of $\Delta\tau=\tau_{\mathrm{CPR}}-\tau_{\lambda=1}$ contains zero or lies above it) and should be better than C-PR ($\lambda=0$) on new transfer senders (the interval of $\tau_{\mathrm{CPR}}-\tau_{\lambda=0}$ lies above zero). The secondary part uses the same window and the matched cohort. Its condition is that C-PR's point estimate lies within the 95\% interval of the better of the two layers on the transfer family and on the allowance family, and that C-PR exceeds both layers on the Sybil-stability family. The paired contrasts of Table~\ref{tab:tau-diff-hybrid} are reported beside this rule as a stricter check; they do not decide it. If its primary part fails, the second half of RQ5, whether C-PR beats a walk on either layer alone, receives a negative answer, and the result is reported with its intervals. The first half, whether C-PR improves on a walk over the transfers alone, is not covered by rule~A; rule~B tests it in W1 (Section~\ref{sub:fresh-holdout}). Once labels have been observed, $\lambda$ is not tuned. When rule~A was fixed, no log of the study had been extracted. On the revised scores it is applied with the same conditions, but after the registered results were known, so its verdict there is not a registered test.

The pipeline also counts two further quantities over the label window: revocations, the owners of a positive allowance at $t_1$ that set it to zero, and a drain heuristic, the owners whose new approval to the spender is followed by a transfer of the same token out of the owner's account, within six hours and of at least half the approved amount, or within one hour when the approval is of $10^{30}$ base units or more. Both count events of the five tokens only. Neither is offered as evidence. Both rise with the number of owners a spender has, so they track use and cannot separate trust from it (Section~\ref{sec:holdout-results}).

\dissertationSubsubheading[sub:trader-holdout]{Trader outcomes in W0}

The spender cohort says nothing about contracts that trade. An exploratory run applied the freeze $t_1$ to the GMX~V2 accounts that are contracts, appear in the freeze-date graph in either layer and close at least three positions between 1~April and 30~June~2026 (Section~\ref{sub:cohorts} gives their number). Its labels are the number of profitable closes, the realized gain, and the shares of closes that were profitable and that were not losses. A close is profitable when it is not a liquidation and its base profit is positive, and a loss when it is not a liquidation and its base profit is negative; the realized gain sums the positive base profits of a contract's closes in USD. Every score of the spender holdout was set against these labels with 400 paired resamples, under no rule, and Chapter~\ref{ch:results} reports the run as exploratory. No liquidation label was computed in this run; the liquidation labels of both windows were first computed under the registration of Section~\ref{sub:fresh-holdout}. With so few contract traders, every interval of this run is wide.

\dissertationSubsubheading[sub:fresh-holdout]{Registered window W1: July to September 2026}

In the registered analysis the W0 holdout also showed something that no rule had asked about. Against its transfer layer walked alone, C-PR ranked the spenders that later gained new approval pairs higher, with an interval above zero, while it was worse on new transfer senders (Appendix~\ref{sec:map-registered}). This was seen only after the labels were in, so it counts as a hypothesis and not as a finding, and it is tested on a window that had not been extracted when the test was fixed. The test was written into a registration document after the W0 results were known. That document and the configuration, in which the status of the window was set to registered, were committed to the public repository at 16:49:32~UTC on 9~October~2026 (commit \texttt{0d3a057}), together with the extraction and evaluation code. Scan~B, the only query that reads logs dated after 30~June~2026, ran after that commit. The extraction script refuses to run unless the registration and the configuration are committed without local changes and the status is registered. It writes the commit and the SHA-256 hashes of both files into an extraction manifest, which records that the extraction started one second after the registration commit (Appendix~\ref{ch:appendix-eval}).

The method is the one fixed in the plan; in the revised analysis the amounts, the tokens and the cohort are those of Section~\ref{sec:amount-revision}, and everything else is unchanged. C-PR is used at $\lambda=0.5$, together with the two sensitivity values and S-PR. The solver settings, the SQL, the extraction filter of the matched cohort, the GMX decoder and the paired bootstrap (400 resamples, seed 42) are those of W0. Scores are frozen at $T_{\mathrm{obs}}=$ 30~June~2026, 23:59:59~UTC, so they are the scores of the same-window evaluation, built from every event since 1~October~2023 with ages up to $1{,}004$ days. Labels are counted from 1~July to 30~September~2026.

Two cohorts are evaluated, and Section~\ref{sub:cohorts} gives their sizes. The spender cohort is defined as in W0, as every contract spender holding at least one positive latest allowance at the freeze. The trader cohort is the GMX~V2 accounts that are contracts, are nodes of the freeze-date graph in either layer and close at least three positions in the label window. The spenders keep their two labels. The traders get the liquidation-free close rate, one minus the share of a contract's July--September closes that ended in liquidation. It is the nearest thing to a default that these data contain, it is built from neither graph, and because it is a rate it does not reward a contract merely for trading more. When it was chosen, no score had been set against a liquidation label in either window, and the W0 trader run had used only the profit labels of Section~\ref{sub:trader-holdout}. What was known about liquidations came from the decoding step, which had shown how common they are among the closes of contract accounts in the observation window (Section~\ref{sub:cohorts}).

Rule~B compares C-PR with its transfer layer walked alone, C-PR ($\lambda=0$), through the paired difference $\Delta\tau=\tau_{\mathrm{CPR}}-\tau_{\lambda=0}$:
\begin{itemize}
\item F1, spenders, new approval pairs: the 95\% interval of $\Delta\tau$ lies above zero;
\item F2, spenders, new transfer senders: the lower end of the interval lies above $-\delta$;
\item F3, traders, liquidation-free close rate: the lower end of the interval lies above $-\delta$.
\end{itemize}
C-PR counts as improving on a walk over the transfers alone only if all three hold, so no correction for multiple testing is needed. The margin $\delta=0.01$ comes from W0: it is the half-width of the paired interval of the corresponding contrast on new transfer senders, C-PR minus C-PR ($\lambda=0$) in the registered analysis (Appendix~\ref{sec:map-registered}), rounded up to two decimals, since a smaller difference could not be resolved at these sample sizes. One weakness of F3 was written down with the rule. If neither score predicts liquidation, F3 can pass with both coefficients near zero, so the two coefficients are reported next to the verdict. Three further results are reported in a fixed order and do not enter the decision: C-PR against C-PR ($\lambda=0$) on the liquidation-free close rate as a superiority test (F4), C-PR against C-PR ($\lambda=1$) on the same label (F5), and the primary part of rule~A, rerun on the new window (F6).

Two sensitivity analyses were registered in the same commit. Both are reported next to the decision, and neither enters it. The first repeats F1--F4 with AWP (Section~\ref{sub:awp-graph}) in place of C-PR ($\lambda=0$). The second repeats F1--F6 with the cohort contracts that a graph leaves out kept as isolated nodes instead of scored zero (Section~\ref{sec:rank-divergence}). Only the two layers walked alone change under the second analysis, since every spender is in the allowance graph and every trader in the union graph that C-PR walks. The run that judged rule~B also set EndorseRank against the in-approve degree, AWP against the transfer in-degree and C-PR against EndorseRank. These contrasts were not registered; they were computed after the W1 labels were known and are reported apart from the decision, as are the supplementary contrasts computed later.

\dissertationSubsubheading[sub:neutral-label-protocol]{Outcomes built from neither edge type}

A comparison of the two edge types on the GMX trader labels, which are built from neither, was registered in the same commit as rule~B, before scan~B. As in a preregistration \parencite{nosek2018preregistration}, its contrasts, decision rules and the wording of each conclusion were written down before the data that decide them were extracted. It reuses the trader cohorts and the label code of Sections~\ref{sub:trader-holdout} and~\ref{sub:fresh-holdout} in two windows: W1 (primary; freeze on 30~June~2026, labels from July to September) and W0 (secondary; freeze on 31~March~2026, labels from April to June, the cohort defined in the same way). When the comparison was registered, the W0 liquidation labels had not been computed for any score; the W0 profit labels and their single-score coefficients were known from the exploratory run.

Four pairs set allowance against transfer: in-approve degree and transfer in-degree; EndorseRank with AWP's activity restarts and AWP, the like-for-like PageRank pair; EndorseRank and AWP; and the amount-weighted layers C-PR ($\lambda=1$) and C-PR ($\lambda=0$). C-PR ($\lambda=0$) minus AWP isolates AWP's weights and restarts. All six trader labels are reported: the liquidation-free close rate (primary), a no-liquidation flag, profitable closes, realized gain, and the profitable and non-loss shares of closes. Two contrasts on the liquidation-free close rate in W1 decide: P, in-approve degree minus transfer in-degree, and Q, EndorseRank with activity restarts minus AWP, each with a Bonferroni 97.5\% interval from $B=2{,}000$ paired resamples (seed 42). Because the chance of at least one liquidation rises with the number of closes, $\tau_b$ is also computed within strata of 3--4, 5--9, 10--24 and 25 or more closes in the label window and weighted by each stratum's contract pairs. Recipients, the traders with an in-approve degree above zero at the freeze, are compared with the other traders on every label, and an influence check leaves out the ten recipients of highest in-approve degree.

R1, a count-level lead on liquidation avoidance, needs the 97.5\% interval of P and the 95\% interval of the stratified P above zero in W1, and a positive P in W0. R2, a PageRank-level lead, needs the 97.5\% interval of Q above zero in W1 and a positive Q in W0. R3 applies if P on the profitable or the non-loss share in W1 is below zero with a 95\% interval excluding zero; the allowance signal is then read as specific to liquidation avoidance. The wording of the conclusion for each outcome was fixed in the same commit. If R1 holds, the approval count at the freeze ranked contract traders above the transfer count on later liquidation avoidance, and if it fails there is no count-level lead of allowance over transfer on that outcome. R2 is worded in the same way for the PageRank pair, and its failure reads as the PageRank versions not being resolved. If R3 applies, the signal concerns forced-liquidation avoidance and not trading success. On the revised scores the same contrasts, rules and wording are applied after the registered results were known.

\dissertationSubheading[sec:proxy-evaluation]{Construct-check design}

Every score enters the same-window evaluation on the matched cohort: EndorseRank with both restart rules, AWP, C-PR at its five values of $\lambda$ and S-PR. For each score $M$ and check $q$ the pipeline reports Spearman's $\rho$ and Kendall's $\tau$ between $\mathbf{s}_M$ and $\mathbf{p}_q$ with 95\% bootstrap intervals (Section~\ref{sec:statistical-tests}), each check oriented so that a higher value means higher expected reputation, as in the evaluation of AWP by its authors \parencite{do2023awp}.

The same-window evaluation is organized around two contemporaneous families:
\begin{enumerate}
\item Transfer (the proxies against which \textcite{do2023awp} evaluated AWP): the contract's inbound transfer activity.
\item Allowance: inbound approval activity, where the contract is the spender.
\end{enumerate}
Both families are graph-adjacent: their inputs are the event classes that make up $G_{\text{transfer}}$ and $G_{\text{approve}}$, although each check is a local statistic and not a global score. Strong alignment of AWP with the transfer checks, or of EndorseRank with the allowance checks, is an intended-construct check: it shows coherence within a construct, part of it is built in, and it is not evidence of independent predictive power. A third family, the Sybil-adjusted activity summaries of Section~\ref{sub:proxy-sybil}, enters the secondary part of rule~A, but no hypothesis rests on it.

The temporal holdouts of Section~\ref{sec:holdout-protocol} provide the out-of-window checks. \textcite{packin2024credit} warn that decentralized credit scores can become opaque; publishing explicit formulas (Section~\ref{sec:proxy-formulas}) and keeping same-window checks apart from the time splits reduce that opacity.

\dissertationSubheading[sec:proxy-formulas]{Operational definitions of proxies}

In what follows, $\mathcal{W}$ denotes the matched cohort, and the proxies are computed over the observation window up to $T_{\mathrm{obs}}$. The notation is that of the proxy queries and the evaluation code in the open pipeline.

\dissertationSubsubheading[sub:proxy-transfer]{Transfer family}

Write $\mathcal{T}_{\mathrm{in}}(w)$ for the set of transfers received by contract $w$ in the window and $u_e$ for the sender of transfer $e$. Self-transfers belong to $\mathcal{T}_{\mathrm{in}}(w)$, so a contract that sends tokens to itself counts itself as one sender, and the amounts $x_e$ are summed in US dollars over the five tokens. Define
\begin{align}
\mathrm{in\_degree}(w)
&=
\bigl|\{\,u_e : e\in\mathcal{T}_{\mathrm{in}}(w)\,\}\bigr|,
\label{eq:in-degree}\\
\mathrm{in\_value}(w)
&=
\sum_{e\in\mathcal{T}_{\mathrm{in}}(w)} x_e.
\label{eq:in-value}
\end{align}
Both are baselines defined in the AWP paper \parencite{do2023awp}, on the premise that economically important addresses receive transfers from more distinct counterparties and in greater total value. Unlike PageRank they are local and ignore the senders' reputations, so comparing AWP with them shows how far propagation, recency weighting and restarts at active senders move its ordering away from the raw inflow.

\dissertationSubsubheading[sub:proxy-allowance]{Allowance family}

On the allowance side, $w$ is the spender (owner $\to w$). With $a^{\star}$ the latest allowances of Equation~\eqref{eq:approve-weight}, define
\begin{align}
\mathrm{in\_approve\_degree}(w)
&=
\bigl|\{\,u : a^{\star}(u,w,\cdot)>0\,\}\bigr|,
\label{eq:in-approve-degree}\\
\mathrm{in\_approve\_value}(w)
&=
\sum_{u,\mathrm{token}} \min\bigl(a^{\star}(u,w,\mathrm{token}),U\bigr).
\label{eq:in-approve-value}
\end{align}
Like the transfer baselines, they count and sum incoming approvals without following multi-hop endorsement chains; the in-approve value uses the capped amounts of the allowance layer of C-PR.

\dissertationSubsubheading[sub:proxy-sybil]{Sybil-adjusted stability family}

The inbound counters leave out self-transfers ($u=w$). Let $N_{\mathrm{tx}}(w)$ be the number of transfers that $w$ receives from other addresses and $N_{\mathrm{uniq}}(w)$ the number of distinct inbound counterparties. Tenure and active months use a second set of rows. Every transfer that involves a cohort contract is attributed to its sender if the sender is in the cohort, and otherwise to its recipient; a transfer between two cohort contracts thus counts only for the sender. We take $t_{\min}(w)$ and $t_{\max}(w)$ to be the earliest and latest timestamps among the rows attributed to $w$. $N_{\mathrm{mon}}(w)$ is the larger of two counts of distinct calendar months: the months in which $w$ received a transfer from another address, and the months of the rows attributed to $w$. Define
\begin{align}
\mathrm{inbound\_counterparty\_ratio}(w)
&=
\frac{N_{\mathrm{uniq}}(w)}{\max(N_{\mathrm{tx}}(w),1)},
\label{eq:counterparty-ratio}\\
\mathrm{transfer\_tenure\_days}(w)
&=
\frac{t_{\max}(w)-t_{\min}(w)}{86400}\quad\text{(seconds to days)},
\label{eq:tenure}\\
\mathrm{active\_months}(w)
&=
N_{\mathrm{mon}}(w).
\label{eq:active-months}
\end{align}
The counterparty ratio penalizes contracts that receive many transfers from few addresses, a simple wash pattern. It does not penalize the usual Sybil pattern, in which many fresh addresses send one transfer each; a contract fed that way scores close to the maximum of one. Tenure and active months track how persistent on-chain activity is over the observation window. These proxies are therefore weak Sybil checks, and colluding adversaries can satisfy all three. They rest on an assumption of this study: activity that lasts and involves many counterparties says more about reputation than short bursts of transfers from few addresses.

\dissertationSubheading[sec:statistical-tests]{Statistical tests}

Neither the reputation scores nor the proxies are binary default indicators; both are continuous and are handled as ranks. As in the evaluation of AWP \parencite{do2023awp}, concordance is measured with rank correlations, which are invariant under any monotonic re-scaling of either variable. Each coefficient is reported with a bootstrap confidence interval.

\begin{itemize}
\item Spearman's $\rho$: the Pearson correlation of the mid-ranks $R_i$ and $S_i$ of the raw score and the raw proxy value of contract $i$, tied values sharing the average of their positions (Section~\ref{sec:rank-correlation-theory}). It measures monotonic agreement: if higher reputation consistently goes with higher proxy values, $\rho$ is positive whatever the nonlinearity in the raw scores.

\item Kendall's $\tau$: every reported value is the tie-corrected $\tau_b$ of Equation~\eqref{eq:kendall}, computed on the raw scores and proxies. Without ties, $(1+\tau)/2$ is the probability that a randomly chosen pair is ordered the same way by both rankings. With ties that reading belongs to $\tau_a=(n_c-n_d)/n_0$ and not to $\tau_b$. Ties are rare among the PageRank scores of the matched cohort (Table~\ref{tab:tie-sensitivity}), but they are common in the raw degrees, which are integer counts, in the holdout labels, which are zero for most spenders, and in the trader labels, whose rates are often one.

\item Why both: the AWP protocol \parencite{do2023awp} reports both, and reporting both keeps the conclusions from hinging on one coefficient. The score distributions are heavy-tailed, with most contracts near zero and a few hubs holding much of the mass, which is why rank coefficients are used instead of a Pearson correlation on raw scores, which the hubs would dominate.

\item Family aggregation: A family is summarized by the plain average over its proxies. The transfer-family mean $\tau$, for example, is $\tfrac{1}{2}(\tau_{\mathrm{in\_degree}}+\tau_{\mathrm{in\_value}})$. On a fixed cohort and fixed proxies, a comparison between scores is only a relative statement about those scores and implies no universal property of reputation. Our conclusions are drawn from the confidence intervals, and no absolute $\tau$ cutoff is used as a decision criterion.

\item Uncertainty: Sampling uncertainty is estimated with a paired bootstrap over contracts \parencite{efron1979bootstrap}. The evaluated cohort (the matched cohort in the same window, the spender or trader cohort in a holdout) is resampled with replacement $B=400$ times (seed 42, fixed in the configuration), and one index matrix is reused for every method and proxy or label, so the resampled statistics are paired across methods. For each $\tau$, family mean and inter-method difference, the 2.5th and 97.5th percentiles of the bootstrap distribution form a 95\% percentile interval \parencite{diciccio1996bootstrap}. Because $\tau_b$ is a smooth function of U-statistics, the percentile interval is first-order accurate at the size of the matched cohort and of the spender cohorts; for the trader cohorts, which are small (Section~\ref{sub:cohorts}), the intervals are wide and their coverage is less certain. The small number of resamples keeps the $O(B\,n\log n)$ cost of the tie-corrected $\tau$ tractable across all methods and proxies. Its price is a Monte Carlo error in each interval end, which for a normal-shaped bootstrap distribution is about $0.13$ times the bootstrap standard error. Only the analysis of Section~\ref{sub:neutral-label-protocol} departs from this, with $B=2{,}000$.

\item Difference tests: Every comparison is a paired contrast on the same resamples. Resample by resample we compute $\Delta\tau=\tau_A-\tau_B$ and report the 95\% percentile interval of its bootstrap distribution; a contrast counts as a difference when its interval excludes zero. Because resampling is paired, the contrast keeps the covariance of the two estimates, $\operatorname{Var}(\hat\tau_A-\hat\tau_B)=\operatorname{Var}\hat\tau_A+\operatorname{Var}\hat\tau_B-2\operatorname{Cov}(\hat\tau_A,\hat\tau_B)$, and with positive covariance its interval is narrower than the two separate intervals suggest. The contrasts that decide a rule were fixed before the data they are judged on: those of rule~A (Section~\ref{sec:holdout-protocol}) before any log was extracted, and those of rule~B and the decision contrasts P and Q (Sections~\ref{sub:fresh-holdout} and~\ref{sub:neutral-label-protocol}) before the W1 logs were extracted, together with F4--F6 and the two sensitivity analyses, which are reported beside the decision. These dates hold for the registered analysis; on the revised scores the same contrasts are computed after the registered results were known (Section~\ref{sec:amount-revision}). C-PR against its transfer layer walked alone on new approval pairs in W0 was not among them; it is marked as exploratory wherever it appears. In the same window, eight contrasts set C-PR and S-PR against whichever layer, walked alone, leads each family (the transfer walk on transfer, the allowance walk on allowance, both on Sybil stability), and four set C-PR against AWP on transfer, against EndorseRank on allowance and against both on Sybil stability. None of these twelve was fixed in the plan; they were computed with the same-window results and are reported beside the secondary part of rule~A. Five same-window contrasts that involve EndorseRank and AWP, none fixed in advance either, are reported in Appendix~\ref{ch:appendix-er-awp}.

\item Directionality and imputation: Each proxy is oriented so that a higher value means a better reputation. After re-alignment to the evaluated cohort, a contract absent from a given proxy source is imputed as zero. With sparse transactions this is a conservative choice, and both scoring approaches apply it in the same way. After this step every evaluated contract has a value in every score and proxy vector, so each correlation uses the whole cohort.

\item Non-inferiority: Two of the three conditions of rule~B ask only that C-PR not fall behind its transfer layer walked alone by more than a margin fixed in advance, $\delta=0.01$. Such a contrast passes when the lower end of its 95\% interval lies above $-\delta$. The test is weaker than a difference test and says nothing about superiority on that label. It is used where the claim under test is that C-PR gives up nothing that matters.

\item What is excluded: The study attempts no causal inference, trains no supervised models on separate training and testing sets, and reports no $p$-values; it asks whether the interval of a contrast excludes the null value.

\end{itemize}

\dissertationSubheading[sec:computational-benchmark]{Computational benchmark (Claim~C1)}

Claim~C1 concerns the computational cost of each score. The benchmark times the solver call that turns an edge list into scores. It leaves out the one-time cost of aggregating and loading the pair tables, which a scoring service that refreshes its scores periodically would amortize.

\begin{itemize}
\item Timed component: The edge arrays and AWP's restart weights are built before the timer starts. The timed call indexes the nodes, assembles the sparse transposed transition matrix, forms the restart vector and runs power iteration with $d=0.85$, $\varepsilon=10^{-8}$ and $I_{\max}=300$. Matrix assembly is vectorised: node identifiers are mapped through a lookup array, the out-strengths come from one weighted count over the source nodes, and the matrix is built in one sparse conversion. Assembly and each iteration therefore both take time linear in $|E|$, and a deployment would repeat both at every update. The calls timed are EndorseRank, AWP, C-PR at $\lambda=0.5$ and S-PR, the last including the solve of the allowance layer that seeds it. The same graphs give the alignment scores.

\item Measurement protocol: Every configuration is measured in one session on an otherwise idle machine. An untimed warm-up run, on which alone peak memory is traced, is followed by five timed runs; we report the mean and standard deviation of the five run times, the number of iterations, and the nodes $|V|$ and edges $|E|$ of the solver graph. Peak memory is the largest allocation traced by Python's \texttt{tracemalloc} module during the warm-up call. The same graph under the same settings occurs in the main benchmark, in the $d=0.85$ rows of the damping sweep and in the scaling stage with $n=15{,}107$, which is the full cohort graph. Those rows reuse the one measurement, and the tables say so; they are not independent measurements of one quantity. Every runtime table lists node and edge counts, since edges drive the cost of an algorithm with $O(|E|)$ work per iteration \parencite{leskovec2020mining}, and the runtimes are written to the machine-readable output (Section~\ref{sec:reproducibility}).

\item Cohorts: The main benchmark runs on the graphs of the matched cohort at $T_{\mathrm{obs}}$, the graphs of the same-window evaluation. For the growth of runtime with graph size we draw deterministic subsamples of the cohort (Section~\ref{sec:scaling-robustness}).

\item Fairness: Every algorithm runs with the same solver, tolerance, iteration cap, damping (outside the damping sweep) and graph construction, so differences in cost come from the graphs and from the restart vectors, which can change the number of iterations, and not from the implementation. The benchmark summary records the platform, the processor and the Python version, Appendix~\ref{sec:pipeline-stages} names the machine and the library versions, and the cost of two scores is compared through time ratios, each reported beside the corresponding edge ratio.

\end{itemize}

\dissertationSubheading[sec:scaling-robustness]{Scaling and robustness extensions}

We also examine how runtime grows with the size of the graph, and how far the family-level pattern moves when the damping, the slope of the value weight or the evaluation sample changes. These analyses qualify Claim~C1 and test the reach of the family-level conclusions; none of them supplies an alternative primary sample. Unless stated otherwise, alignment claims in Chapter~\ref{ch:results} refer to the matched cohort.

\dissertationSubsubheading[sub:expanded-pool]{Cohort subsamples for runtime scaling}

The scaling stages are deterministic subsamples of the matched cohort. Contracts are sorted by
\begin{center}
\texttt{SHA256(}\textit{seed}\texttt{:}\textit{address}\texttt{)} with \textit{seed} $=$ \texttt{contract-benchmark-v1},
\end{center}
that is, by the hash of one string made of the seed, a colon and the lower-cased address, and the stage of size $n$ takes the first $n$ contracts in that order. The seed string is a constant fixed in the configuration of 9~October~2026. The stages are $n=1{,}000$, $2{,}000$, $5{,}000$, $10{,}000$ and the full cohort, $n=15{,}107$. Since the hashing is deterministic, the subsamples cannot drift between library versions the way pseudorandom draws do, and each stage contains the one before it.

At every stage the graph keeps the edges of the cohort graph at $T_{\mathrm{obs}}$ that have at least one end point among the sampled contracts, and the solver's node set consists of the end points of those edges. The edge count therefore grows with the stage, up to the full cohort graph, but not in proportion to $n$: a contract with many counterparties brings many edges, and an edge between two sampled contracts is counted once. The stage graphs are ego networks of subsets of one cohort, not samples of a larger chain, and no stage goes beyond the full cohort graph. The experiment therefore shows how runtime grows with $|E|$ up to that size and says nothing about larger graphs. Every scaling table reports per-stage $|V|$ and $|E|$ for both methods, and the scaling statement in Chapter~\ref{ch:results} goes no further than those columns allow.

\dissertationSubsubheading[sub:damping-protocol]{Damping sensitivity}

On the matched cohort we vary the damping over $d\in\{0.75,0.85,0.95\}$, the values fixed in the plan, with the graphs and proxies held fixed. A smaller $d$ raises teleportation; a larger one gives multi-hop paths more weight and may slow convergence. The sweep asks whether alignment on the allowance and transfer families depends on the teleportation probability.

\dissertationSubsubheading[sub:top-tokens]{Value-slope sensitivity}

The plan fixed a check on subgraphs of twenty tokens, chosen by summed raw amount and by number of rows. Once every graph keeps five tokens that check loses its point, and it is replaced by a check of the parameter that the revision introduced (Section~\ref{sec:amount-revision}). EndorseRank and AWP are recomputed at $T_{\mathrm{obs}}$ with the slope $b$ of $V$ divided and multiplied by ten, which moves the point of half weight to ten times and to a tenth of the median amount. The check asks whether the family-level results depend on how steeply the amount weight rises. The registered analysis keeps its token check, published with its results.

\dissertationSubsubheading[sub:sample-size-protocol]{Sample-definition sensitivity}

At each scaling stage of Section~\ref{sub:expanded-pool} the family means of EndorseRank and AWP are recomputed on the stage's contracts, scored on the stage's graph. A stage differs from the full cohort in size and in its graph: a contract's score changes because the edges of the contracts left out of the stage are missing. The per-stage coefficients are therefore a sensitivity check on the sample definition and not evidence of stability, and no intervals were computed for them; the question is how far each coefficient moves when the sample changes. This analysis was not in the plan; it was run with the robustness checks after the W1 labels were known.

\dissertationSubheading[sec:reproducibility]{Reproducibility}

No table in Chapter~\ref{ch:results} was transcribed by hand. Once the Phase~1 tables are on disk, the evaluation drivers recompute every check, every alignment statistic with its interval, every holdout result and the benchmark, and write them to machine-readable summaries. An export script writes the tables and figures of Chapter~\ref{ch:results} and Appendix~\ref{ch:appendix-eval} from those summaries, and a publishing script copies the summaries with their SHA-256 checksums and a digest that names, for every number, the file and the key it comes from. One versioned configuration file holds the windows, the cohort threshold, the event topics, the damping, the decay parameters, the values of $\lambda$, the bootstrap count and seed, the benchmark protocol and seed, the cost budget and the part size. The revised analysis adds the token list, the daily price table and the constants of Section~\ref{sec:amount-revision} as files in the repository, and one switch selects which of the two analyses the scripts compute.

The design was fixed in commits to the public repository at \url{https://github.com/abitocodes/phd_works}: the analysis plan and the configuration on 9~October~2026 (commit \texttt{43d8045}), before any log was extracted; the registration of W1 together with the extraction and evaluation code at 16:49~UTC on the same day (commit \texttt{0d3a057}), before scan~B; and the revision of the amounts on 10~October~2026 (commits \texttt{a897475}, \texttt{09dfc32} and \texttt{2705b28}), before any revised quantity was computed. The reproducible elements are these: deterministic subsampling of the cohort by SHA256 hashes; a fixed PageRank tolerance and iteration cap; one bootstrap seed and one resample-index matrix for each cohort; the proxy formulas of Section~\ref{sec:proxy-formulas}; the benchmarking procedure of Section~\ref{sec:computational-benchmark}, with the time of every run kept in the output; the coupling weight $\lambda$ and the holdout logic of Section~\ref{sec:holdout-protocol}, both fixed in the configuration before any log was extracted; and the extraction manifest of scan~B with the hashes of the registered files. From the downloaded tables the alignment, holdout and benchmark results can be rebuilt without any new BigQuery query. The machine-readable summaries behind all the tables, together with the configuration, the analysis plan and the registration, are committed to the repository under \texttt{experiments/data/3-published-results-for-thesis/}, those of the revised analysis in its subfolder \texttt{revised-usd/}, in the same project folder as the source code. The numbers in this thesis can thus be checked against those summaries without rerunning the pipeline. Each summary records the code revision that wrote it, and the digest lists any script changed since that revision. Appendix~\ref{ch:appendix-eval} lists the repository URL, the commit hashes, the configuration file and the summary files.

\dissertationSubheading[sec:sampling-bias]{Sampling considerations and missingness}

The matched cohort is a constructed sample, not a random draw from all Arbitrum contracts: eligibility needs non-zero approvals of one of the five tokens from at least three distinct owners in the window, and transfer activity plays no part in it. Same-window findings therefore apply to contracts on Arbitrum One that at least three owners approved on WETH, USDC, USDT0, ARB or WBTC in these months, and holdout findings to the contract spenders in the cohort's approval data. Contracts that fewer than three owners approved, and contracts that no owner approved at all, are excluded by design, and cohort contracts outside a graph score $0$ under that method (Section~\ref{sec:data-sources}). Because the cohort is defined from approvals, it favors the allowance side in coverage: every cohort contract received approvals at some time in the window, while AWP scores zero the cohort contracts that received no transfer (Table~\ref{tab:tie-sensitivity}).

The spender cohorts of the holdouts include spenders outside the matched cohort, which entered the extract because a cohort contract approved them. For such a spender both the scores and the labels see only the approvals of the five tokens that cohort contracts granted it and the transfers of those tokens it exchanged with cohort contracts, so its counts understate its activity on the chain. The trader cohorts are small (Section~\ref{sub:cohorts}) because few contract accounts closed GMX~V2 positions in these months, and every trader result has wide intervals.

Data can be missing in a few ways. An amount that cannot be read gives zero (Section~\ref{sec:event-decoding}). EndorseRank does not see approvals that emit no standard ERC-20 event, nor allowances set before 1~October~2023 and not changed since, and AWP does not see transfers made off-chain or inside custodial accounts. Both are thus biased toward on-chain, standard-token behavior, which is acceptable for an on-chain reputation study but matters once results are extrapolated to all contracts of a protocol. The revision adds three kinds of missing data: tokens other than the five, which are left out by design and counted; transfers of a token before its first daily price, which are left out and counted; and days without a price point, which take the previous day's price (Appendix~\ref{sec:data-tokens}).

\dissertationSubheading[sec:implementation-notes]{Implementation notes}

The pair tables, the proxies and the labels are computed in BigQuery SQL. The PageRank code is written in Python, with a SciPy sparse transition matrix, NumPy power iteration and vectorised matrix assembly, so that graphs with tens of millions of edges fit one solver call. Each configuration builds its edge arrays once. The timing benchmarks leave out I/O and the aggregation of the pair tables and cover operator assembly and iteration only (Section~\ref{sec:computational-benchmark}). The benchmark ran on one machine, an Intel Core Ultra~9 185H processor with 64~GB of memory under Windows~11, with Python 3.11.9, NumPy 2.3.2 and SciPy 1.16.1 (Appendix~\ref{sec:pipeline-stages}). Randomness enters only through the bootstrap, whose seed is fixed; given the downloaded tables and the configuration, evaluation is deterministic. Small floating-point differences between platforms are to be expected, but they would be negligible against the three-decimal-place precision of the $\tau$ metric.

Figures are generated from the same evaluation summaries as the tables, so the plotted values agree with the tabulated numbers.

\dissertationSubheading[sec:ethics]{Ethics}

Only publicly available blockchain data are used in this research, and it involves no human subjects. The study was reviewed by the institutional ethics committee (Unisa SoC/CAES). Appendix~\ref{sec:ethics-certificate} contains the ethics clearance certificate (reference [Unisa SoC/CAES ethics clearance ref: TBD], dated [TBD]).

\begin{itemize}
\item\phantomsection\label{sub:privacy-anonymity}
  Privacy and anonymity: A blockchain address is a pseudonym and does not by itself reveal the real-world identity behind it \parencite{packin2024credit}. Pseudonymity is not the same as anonymity, however, and in theory a determined adversary could tie an address to a person through exchange deposits, online profiles or heuristics. This study makes no such links. The scored addresses are contracts; EOAs, which may belong to individuals, appear only as owners, senders and recipients, and no result is reported for any of them. Our data collection is limited to on-chain \texttt{Approval}, \texttt{Transfer} and GMX~V2 position events and to the code stored at an address, and we collected no off-chain personally identifiable information (such as IP addresses, exchange KYC records, social media handles or deanonymization graphs). We did not cluster addresses into real-world entities, publish leaderboards of addresses likely to be personally identifiable, or try to join an address's transaction history with other databases to infer identity. The published results are aggregate measures only (correlation coefficients, runtimes, family means) and contain no individualized rankings. Where the main text quotes examples, they name broad categories (for example, `spenders with high allowances') and never specific addresses.


\item\phantomsection\label{sub:data-minimization}
  Data minimization: Scan~A keeps every ERC-20 \texttt{Approval} and \texttt{Transfer} event of the window in a table of the study's own cloud project, because the cohort rule needs every approval. All later queries and every downloaded table are restricted to the ego network of the cohort and to the fields that graph construction and the checks require, namely the relevant addresses, token amounts and timestamps. The extracted tables are kept so that the research pipeline can be audited. They will not be released for the purpose of deanonymizing users.


\item\phantomsection\label{sub:data-usage-compliance}
  Data usage and compliance: Every query ran against public data hosted in the \texttt{bigquery-public-data.goog\_blockchain\_arbitrum\_one\_us} dataset and complied with the Google Cloud terms of service and rate limits. A dry run before every query and a budget for the whole study prevented unnecessary scans. The account-type check sent read-only \texttt{eth\_getCode} calls to public RPC endpoints, and the token decimals and the Chainlink feed values were read with read-only \texttt{eth\_call} requests; daily token prices came from the public DefiLlama API. We sent no transactions and accessed no privileged or non-public indexing endpoints, and no smart contract was invoked in any way that might interfere with the normal operation of the network.


\item\phantomsection\label{sub:fairness-transparency}
  Fairness and transparency: Decentralized credit scores raise concerns about fairness as well as opacity \parencite{packin2024credit}. Graph-based contract scores can undervalue new contracts, rarely used contracts and contracts that receive no public approvals, and both EndorseRank and AWP carry structural biases of this kind. In the allowance graph, for instance, a contract that has received no spending authority is at a disadvantage relative to contracts that have; transfer graphs, in turn, favor addresses embedded in value-flow networks. Chapter~\ref{ch:discussion} discusses the limitations, among them cold-start effects and the Sybil susceptibility of approval graphs. The code, the parameters and the evaluation metrics are in the public repository (Appendix~\ref{ch:appendix-eval}), so that others can review the work and repeat it.


\item\phantomsection\label{sub:responsible-disclosure}
  Responsible disclosure and impact: EndorseRank may inform risk assessment in DeFi, but it is not financial advice, credit advice, a security rating or a consumer reporting product. Lending, trading and compliance decisions still need human oversight. Chapter~\ref{ch:discussion} discusses the potential for adversarial manipulation, for example Sybil attacks on $G_{\text{approve}}$ or wash transfers on $G_{\text{transfer}}$. The study reports aggregated, cohort-level findings and names no contract as trustworthy or untrustworthy. Any future operational deployment would require additional governance, appeal mechanisms and fairness monitoring, all of which lie beyond the scope of this research.


\item\phantomsection\label{sub:methodological-commitments}
  Summary of methodological commitments: The analysis plan fixed the chain, the windows, the cohort and account-type rules, the graph formulas, the proxy definitions, the spender labels, the bootstrap and timing protocols, the statistical tests and rule~A before any log was extracted. The registration fixed the trader labels, rule~B and the comparison on trader outcomes after the W0 results were known and before the W1 logs were extracted. The revision of 10~October~2026 changed the amounts, the token set and with them the cohort, after the registered results were known and before any revised quantity was computed (Section~\ref{sec:amount-revision}). The primary claims concern all the scores together, on the matched contract cohort and the spender holdouts, with each PageRank set against the raw degree it smooths. The cohort subsamples show how runtime grows with $|E|$, and the sensitivity sweeps show how far results depend on parameters and on the sample definition.
\end{itemize}

Chapter~\ref{ch:results} reports what this protocol yields for Arbitrum One: the dataset and cohorts, the runtime benchmark, the same-window alignment of every score, the rank divergence, the sensitivity checks, the temporal holdouts W0 and W1 with the rules fixed for them, the outcomes built from neither edge type, and finally Claims~C1--C6.
